% =====================================================================
%  Anatomía de un LLM — Fase I — Módulo 2
%  Álgebra lineal I: vectores, matrices y embeddings
%  Compilar: latexmk -r ../latexmkrc M02-algebra-lineal-1.tex
% =====================================================================
\documentclass[11pt,a4paper,twoside,openany]{book}
\usepackage{../estilo-apuntes}

\newcommand{\ApuntesCurso}{Anatomía de un LLM}
\newcommand{\ApuntesModulo}{Módulo 2}

% Notación propia del módulo
\newcommand{\ip}[2]{\langle #1,#2\rangle}
\newcommand{\unos}{\mathbf{1}}
\DeclareMathOperator{\Var}{Var}
\DeclareMathOperator{\diag}{diag}

\hypersetup{pdftitle={Módulo 2. Álgebra lineal I},pdfauthor={Anatomía de un LLM}}

\begin{document}
\frontmatter
% ------------------------------------------------------------------ Portada
\begin{titlepage}
\centering
\vspace*{2.5cm}
{\large\scshape Anatomía de un LLM}\par
\vspace{0.4cm}
{\normalsize Fundamentos matemáticos de los modelos de lenguaje}\par
\vspace{2.2cm}
\rule{\textwidth}{0.8pt}\par\vspace{0.5cm}
{\small\scshape Fase I \textperiodcentered\ Fundamentos matemáticos}\par\vspace{0.3cm}
{\Huge\bfseries Módulo 2\par}\vspace{0.3cm}
{\LARGE Álgebra lineal I\par}
\vspace{0.2cm}
{\large Vectores, matrices y embeddings\par}
\vspace{0.5cm}\rule{\textwidth}{0.4pt}\par
\vspace{1.6cm}
{\large Apuntes del curso}\par
\vfill
{\small Edición de octubre de 2026}\par
\end{titlepage}

% ------------------------------------------------------- Presentación
\chapter*{Presentación del módulo}
\addcontentsline{toc}{chapter}{Presentación del módulo}
\markboth{Presentación del módulo}{}

Un gran modelo de lenguaje no manipula palabras, sino vectores: cada token se representa por un elemento de $\R^{d}$, y casi todo el cálculo posterior consiste en productos escalares y productos de matrices. Este módulo desarrolla la parte del álgebra lineal que se necesita para entender esa representación: la geometría del producto escalar, las matrices como aplicaciones lineales y la composición de aplicaciones mediante el producto matricial. Todo ello se aplica después a la capa de embeddings, a las analogías entre palabras y, como adelanto, a la puntuación de la atención.

\section*{Objetivos}
Al finalizar el módulo, el estudiante debe ser capaz de:
\begin{enumerate}[label=\arabic*.]
  \item operar con vectores de $\R^{d}$ y enunciar las propiedades del producto escalar y de las normas;
  \item demostrar la desigualdad de Cauchy--Schwarz y la desigualdad triangular, y usarlas para definir el ángulo entre vectores en cualquier dimensión;
  \item calcular proyecciones sobre una recta y similitudes coseno, y relacionar la similitud coseno con la distancia euclídea;
  \item interpretar una matriz como una aplicación lineal y demostrar que toda aplicación lineal entre espacios $\R^{n}$ es de esa forma;
  \item calcular productos de matrices de varias maneras, justificar sus propiedades y estimar su coste;
  \item describir la capa de embeddings como una matriz $K\times d$, resolver analogías con aritmética de vectores y explicar el escalado por $\sqrt{d}$ de la atención.
\end{enumerate}

\section*{Requisitos previos}
Trigonometría elemental (seno, coseno y teorema del coseno), resolución de sistemas $2\times 2$ y manejo de sumatorios. Del módulo~1 se usan el vocabulario $\V$ y los índices de los tokens (tema~1), la matriz de transición de un modelo de bigramas (tema~2) y la función softmax (tema~4). En el tema~2 y en el ejercicio~\ref{eje:5-atencion} se emplean la esperanza y la varianza de variables aleatorias y su comportamiento con la independencia, que se estudian con detalle en el módulo~5.

\section*{Organización}
\begin{center}
\begin{tabular}{@{}cll@{}}
\toprule
Tema & Contenido & Fuentes principales \\
\midrule
1 & Vectores, producto escalar y norma & \cite{strang}, \cite{mml}, \cite{vmls} \\
2 & Ángulos, ortogonalidad y similitud coseno & \cite{strang}, \cite{mml}, \cite{bhk} \\
3 & Matrices y aplicaciones lineales & \cite{strang}, \cite{mit1806sc}, \cite{mml} \\
4 & El producto de matrices & \cite{strang}, \cite{mit1806sc}, \cite{vmls} \\
5 & Embeddings, analogías y atención & \cite{jm3}, \cite{mikolov-a}, \cite{vaswani} \\
\bottomrule
\end{tabular}
\end{center}

Cada tema termina con una lista de lecturas y con ejercicios graduados en tres niveles: $\star$ comprensión, $\star\star$ cálculo y $\star\star\star$ demostración. Los ejercicios marcados con \emph{Fuente} están adaptados de cursos y libros de texto universitarios; el enunciado se ha redactado de nuevo y, en algunos casos, se han cambiado los datos para que puedan resolverse a mano. Las soluciones completas, paso a paso, están en el \cref{ap:soluciones}.

\section*{Sobre las fuentes}
El texto es de redacción propia. La secuencia de los temas~1 a~4 sigue los capítulos~1 y~2 del libro de Strang~\cite{strang}, texto de los cursos 18.06 y 18.06SC del MIT~\cite{mit1806,mit1806sc}, y los capítulos~2 y~3 de Deisenroth, Faisal y Ong~\cite{mml}, de donde se toman las definiciones abstractas de norma, producto interior y aplicación lineal. El libro de Boyd y Vandenberghe~\cite{vmls}, texto del curso EE103 de Stanford, aporta el punto de vista computacional (coste de las operaciones, vectores de recuento de palabras). La geometría en alta dimensión sigue el capítulo~2 de Blum, Hopcroft y Kannan~\cite{bhk}. El tema~5 se basa en el capítulo~5 de Jurafsky y Martin~\cite{jm3}, en la clase~2 del curso CS224N~\cite{cs224n} y en los artículos originales citados en él. Sobre los modelos comerciales (Claude, Gemini, Grok) solo se usa información pública.

\section*{Notación}
\begin{center}
\begin{tabular}{@{}ll@{}}
\toprule
Símbolo & Significado \\
\midrule
$u=(u_1,\dots,u_d)\in\R^{d}$ & vector columna de $d$ componentes \\
$e_1,\dots,e_d$ & vectores de la base canónica (o vectores \emph{one-hot}) \\
$u\cdot v=u^{\top}v$ & producto escalar \\
$\norm{u}$, $\norm{u}_1$, $\norm{u}_\infty$ & normas euclídea, $\ell_1$ y del máximo \\
$\hat u=u/\norm{u}$ & vector normalizado ($u\neq 0$) \\
$\cos(u,v)$ & similitud coseno \\
$A\in\R^{m\times n}$, $A_{ij}$ & matriz de $m$ filas y $n$ columnas y su entrada $(i,j)$ \\
$A^{\top}$, $I_n$ & traspuesta e identidad de orden $n$ \\
$\V$, $K=\abs{\V}$, $E\in\R^{K\times d}$ & vocabulario, su tamaño y matriz de embeddings \\
$q$, $k$, $d_k$ & consulta, clave y su dimensión (atención) \\
\bottomrule
\end{tabular}
\end{center}

Los vectores se escriben con letras minúsculas y se consideran columnas; cuando se escriben en línea, $(u_1,\dots,u_d)$ denota la columna correspondiente. Los índices de los tokens empiezan en~0, como en el módulo~1; los de las componentes de un vector, en~1.

\tableofcontents

\mainmatter
% =====================================================================
\chapter{Vectores, producto escalar y norma}
\label{tema:1}
% =====================================================================

Este tema introduce el espacio $\R^{d}$, sus dos operaciones básicas y las dos herramientas métricas sobre las que descansa todo el módulo: el producto escalar y la norma. La presentación sigue las secciones~1.1 y~1.2 de Strang~\cite{strang} y la clase~1 del curso 18.06SC~\cite{mit1806sc}; las definiciones abstractas de espacio vectorial, norma y producto interior proceden de~\cite[§2.4, §3.1 y §3.2]{mml}, y el punto de vista computacional, de~\cite[cap.~1 y §3.1]{vmls}.

\section{El espacio \texorpdfstring{$\R^{d}$}{ℝᵈ}}

\begin{definicion}[Vector]
Sea $d\ge 1$ un entero. Un \emph{vector} de $\R^{d}$ es una lista ordenada $u=(u_1,\dots,u_d)$ de $d$ números reales, llamados \emph{componentes} de $u$. Dos vectores son iguales si lo son todas sus componentes. El \emph{vector nulo} es $0=(0,\dots,0)$.
\end{definicion}

Un vector de $\R^{2}$ o de $\R^{3}$ admite dos lecturas geométricas equivalentes: un punto del plano o del espacio, o una flecha que parte del origen y termina en ese punto. Para $d>3$ no hay dibujo posible, pero todas las definiciones de este módulo se formulan componente a componente y valen igual en cualquier dimensión.

\begin{definicion}[Suma y producto por un escalar]
Para $u,v\in\R^{d}$ y $c\in\R$ se definen
\[
u+v=(u_1+v_1,\dots,u_d+v_d),\qquad c\,u=(c\,u_1,\dots,c\,u_d).
\]
\end{definicion}

Geométricamente, $u+v$ se obtiene colocando la flecha $v$ a continuación de la flecha $u$ (regla del paralelogramo), y $c\,u$ estira la flecha si $c>1$, la encoge si $0<c<1$ y la invierte si $c<0$. La diferencia $u-v=u+(-1)v$ es la flecha que va del extremo de $v$ al extremo de $u$.

\begin{proposicion}\label{prop:ev}
Con la suma y el producto por un escalar, $\R^{d}$ satisface, para todos $u,v,w\in\R^{d}$ y $a,b\in\R$:
\begin{enumerate}
  \item $u+v=v+u$ y $(u+v)+w=u+(v+w)$;
  \item $u+0=u$ y $u+(-u)=0$;
  \item $a(u+v)=au+av$, $(a+b)u=au+bu$ y $(ab)u=a(bu)$;
  \item $1\,u=u$.
\end{enumerate}
\end{proposicion}

\begin{proof}
Cada igualdad se reduce, componente a componente, a la propiedad correspondiente de los números reales. Por ejemplo, la componente $i$ de $a(u+v)$ es $a(u_i+v_i)=au_i+av_i$, que es la componente $i$ de $au+av$.
\end{proof}

Un conjunto dotado de una suma y de un producto por escalares reales que cumplen las propiedades de la \cref{prop:ev} se llama \emph{espacio vectorial real}~\cite[§2.4]{mml}. Además de $\R^{d}$, son espacios vectoriales el conjunto $\R^{m\times n}$ de matrices (\cref{tema:3}) y el de polinomios de grado a lo sumo $n$. Este módulo trabaja casi siempre en $\R^{d}$.

\begin{definicion}[Combinación lineal]
Sean $v_1,\dots,v_k\in\R^{d}$ y $c_1,\dots,c_k\in\R$. El vector $c_1v_1+\dots+c_kv_k$ es una \emph{combinación lineal} de $v_1,\dots,v_k$ con coeficientes $c_1,\dots,c_k$.
\end{definicion}

\begin{ejemplo}\label{ej:operaciones}
Sean $u=(1,2)$ y $v=(3,-1)$. Entonces $u+v=(4,1)$, $2u=(2,4)$ y $u-v=(-2,3)$. La combinación lineal $2u+v$ vale $(5,3)$.
\end{ejemplo}

Las combinaciones lineales son la operación central del álgebra lineal~\cite[§1.1]{strang}: el producto de una matriz por un vector (\cref{tema:3}) es una combinación lineal de las columnas de la matriz, y la cuestión de qué vectores se obtienen como combinación de otros dados conduce a los conceptos de independencia, base y rango del módulo~3.

\begin{observacion}[Vectores en los modelos de lenguaje]
Un modelo de lenguaje asigna a cada token un vector de $\R^{d}$, su \emph{embedding} (\cref{tema:5}). En el modelo más pequeño de GPT-2 es $d=768$~\cite{radford}, y en GPT-3, con 175\,000 millones de parámetros, $d=12\,288$~\cite[tabla~2.1]{brown}. Anthropic, Google y xAI no publican este valor para los modelos actuales de Claude, Gemini y Grok. Otro ejemplo clásico es el \emph{vector de recuentos} de un documento: fijado un diccionario de $n$ palabras, la componente $i$ es el número de veces que aparece la palabra $i$~\cite[§1.1]{vmls}.
\end{observacion}

\section{El producto escalar}

\begin{definicion}[Producto escalar]\label{def:escalar}
El \emph{producto escalar} (o producto interior euclídeo) de $u,v\in\R^{d}$ es el número real
\[
u\cdot v=\sum_{i=1}^{d}u_iv_i.
\]
Con la notación matricial del \cref{tema:3}, $u\cdot v=u^{\top}v$.
\end{definicion}

\begin{ejemplo}
$(1,2,3)\cdot(4,5,6)=4+10+18=32$. Obsérvese que el resultado es un número, no un vector: el producto componente a componente $(4,10,18)$ es otra operación.
\end{ejemplo}

\begin{proposicion}[Propiedades del producto escalar]\label{prop:escalar}
Para todos $u,v,w\in\R^{d}$ y $a,b\in\R$:
\begin{enumerate}
  \item (simetría) $u\cdot v=v\cdot u$;
  \item (linealidad en cada argumento) $(au+bv)\cdot w=a\,(u\cdot w)+b\,(v\cdot w)$;
  \item (positividad) $u\cdot u\ge 0$, con igualdad si y solo si $u=0$.
\end{enumerate}
\end{proposicion}

\begin{proof}
(i) Se sigue de $u_iv_i=v_iu_i$. (ii) $\sum_i(au_i+bv_i)w_i=a\sum_iu_iw_i+b\sum_iv_iw_i$. Por la simetría, la linealidad vale también en el segundo argumento. (iii) $u\cdot u=\sum_iu_i^{2}$ es una suma de cuadrados, luego es $\ge 0$, y se anula si y solo si cada $u_i=0$.
\end{proof}

Una aplicación $\ip{\cdot}{\cdot}\colon V\times V\to\R$ sobre un espacio vectorial que cumple (i)--(iii) se llama \emph{producto interior}~\cite[§3.2]{mml}. El producto escalar es el producto interior estándar de $\R^{d}$; todos los resultados de este tema, salvo los que usan componentes, valen para un producto interior cualquiera.

\begin{observacion}[Coste]
Calcular $u\cdot v$ en $\R^{d}$ requiere $d$ multiplicaciones y $d-1$ sumas, unas $2d$ operaciones en coma flotante~\cite[§1.5]{vmls}. Un LLM calcula miles de millones de productos escalares por cada token generado, y su organización en productos de matrices (\cref{tema:4}) es lo que permite ejecutarlos eficientemente en una GPU.
\end{observacion}

\section{Normas}

\begin{definicion}[Norma euclídea]\label{def:norma}
La \emph{norma} (o longitud) euclídea de $u\in\R^{d}$ es
\[
\norm{u}=\sqrt{u\cdot u}=\sqrt{u_1^{2}+\dots+u_d^{2}}.
\]
Un vector $u$ es \emph{unitario} si $\norm{u}=1$. Si $u\neq 0$, el vector $\hat u=u/\norm{u}$ es unitario y tiene la misma dirección y sentido que $u$; se dice que $\hat u$ es el \emph{normalizado} de $u$. La \emph{distancia} entre $u$ y $v$ es $\norm{u-v}$.
\end{definicion}

En $\R^{2}$ y $\R^{3}$ la norma es la longitud de la flecha, por el teorema de Pitágoras. Por ejemplo, $\norm{(3,4)}=5$, $\norm{(1,2,2)}=3$ y el normalizado de $(3,4)$ es $(0{,}6;\,0{,}8)$.

\begin{definicion}[Norma]\label{def:norma-general}
Una \emph{norma} sobre un espacio vectorial $V$ es una aplicación $\norm{\cdot}\colon V\to\R$ que cumple, para todos $u,v\in V$ y $c\in\R$~\cite[§3.1]{mml}:
\begin{enumerate}
  \item (definida positiva) $\norm{u}\ge 0$, con igualdad si y solo si $u=0$;
  \item (homogeneidad) $\norm{c\,u}=\abs{c}\,\norm{u}$;
  \item (desigualdad triangular) $\norm{u+v}\le\norm{u}+\norm{v}$.
\end{enumerate}
\end{definicion}

\begin{ejemplo}[Otras normas de $\R^{d}$]
Además de la euclídea, también llamada norma $\ell_2$, se usan la norma $\ell_1$, $\norm{u}_1=\sum_i\abs{u_i}$, y la norma del máximo, $\norm{u}_\infty=\max_i\abs{u_i}$. Para $u=(3,-4)$: $\norm{u}_1=7$, $\norm{u}=5$ y $\norm{u}_\infty=4$. El cuadrado de la norma euclídea de los parámetros es la penalización de la regularización $L_2$ del módulo~8.
\end{ejemplo}

Las propiedades (i) y (ii) de la \cref{def:norma-general} son inmediatas para la norma euclídea. La desigualdad triangular, en cambio, requiere un resultado previo que es probablemente la desigualdad más importante del álgebra lineal.

\begin{teorema}[Desigualdad de Cauchy--Schwarz]\label{teo:cs}
Para todos $u,v\in\R^{d}$,
\[
\abs{u\cdot v}\le\norm{u}\,\norm{v},
\]
y la igualdad se da si y solo si uno de los dos vectores es múltiplo del otro.
\end{teorema}

\begin{proof}
Si $v=0$, ambos miembros valen~0 y $v=0\cdot u$. Sea $v\neq 0$ y sea $t^{*}=(u\cdot v)/\norm{v}^{2}$. Por la positividad y la linealidad del producto escalar (\cref{prop:escalar}),
\[
0\le\norm{u-t^{*}v}^{2}=\norm{u}^{2}-2t^{*}(u\cdot v)+(t^{*})^{2}\norm{v}^{2}
=\norm{u}^{2}-\frac{(u\cdot v)^{2}}{\norm{v}^{2}}.
\]
Multiplicando por $\norm{v}^{2}>0$ se obtiene $(u\cdot v)^{2}\le\norm{u}^{2}\norm{v}^{2}$, y tomando raíces cuadradas, el enunciado. La igualdad se da si y solo si $\norm{u-t^{*}v}=0$, es decir, si $u=t^{*}v$. Recíprocamente, si $u=cv$, entonces $\abs{u\cdot v}=\abs{c}\norm{v}^{2}=\norm{u}\norm{v}$.
\end{proof}

La elección de $t^{*}$ no es arbitraria: $t^{*}v$ es la proyección de $u$ sobre la recta de $v$ (\cref{prop:proyeccion}), y la demostración dice que el cuadrado de la distancia de $u$ a esa recta es no negativo. El ejercicio~\ref{eje:1-cs} propone una demostración alternativa.

\begin{corolario}[Desigualdad triangular]\label{cor:triangular}
Para todos $u,v\in\R^{d}$, $\norm{u+v}\le\norm{u}+\norm{v}$. En consecuencia, la norma euclídea es una norma en el sentido de la \cref{def:norma-general}.
\end{corolario}

\begin{proof}
Por la linealidad del producto escalar y el \cref{teo:cs},
\[
\norm{u+v}^{2}=\norm{u}^{2}+2\,u\cdot v+\norm{v}^{2}\le\norm{u}^{2}+2\norm{u}\norm{v}+\norm{v}^{2}=\bigl(\norm{u}+\norm{v}\bigr)^{2}.
\]
Ambos miembros son no negativos, luego se puede tomar la raíz cuadrada. Las propiedades (i) y (ii) son inmediatas: $\norm{cu}^{2}=\sum_i c^{2}u_i^{2}=c^{2}\norm{u}^{2}$.
\end{proof}

\begin{observacion}
La desigualdad triangular expresa que el camino directo es el más corto: la longitud del lado $u+v$ de un triángulo no supera la suma de las de los otros dos. Aplicada a $u-w=(u-v)+(v-w)$, da $\norm{u-w}\le\norm{u-v}+\norm{v-w}$, que es la desigualdad triangular para distancias.
\end{observacion}

\begin{lecturas}
\item Strang~\cite{strang}, §1.1 (vectores y combinaciones lineales) y §1.2 (longitudes y productos escalares).
\item MIT 18.06SC~\cite{mit1806sc}, unidad~I, sesión «The Geometry of Linear Equations».
\item Deisenroth, Faisal y Ong~\cite{mml}, §2.4 «Vector Spaces», §3.1 «Norms» y §3.2 «Inner Products».
\item Boyd y Vandenberghe~\cite{vmls}, cap.~1 «Vectors» y §3.1 «Norm».
\end{lecturas}

\section*{Ejercicios del tema 1}
\addcontentsline{toc}{section}{Ejercicios del tema 1}

\begin{ejercicio}{1}{eje:1-combinacion}
Sean $u=(1,2)$ y $v=(3,-1)$. Halle los coeficientes $a,b\in\R$ tales que $au+bv=(5,3)$ e interprete geométricamente el resultado.
\end{ejercicio}

\begin{ejercicio}{1}{eje:1-recuentos}
Sea $w\in\R^{n}$ el vector de recuentos de un documento respecto de un diccionario de $n$ palabras que contiene todas las del documento, y sea $\unos=(1,\dots,1)\in\R^{n}$. (a) ¿Qué representa $\unos\cdot w$? (b) ¿Qué significa que $w_{282}=0$? (c) Se define el \emph{vector histograma} $h$, cuya componente $i$ es la fracción de las palabras del documento que son la palabra $i$. Exprese $h$ en función de $w$. (d) Con el diccionario (el, gato, come, pez), en este orden, calcule $w$, $\unos\cdot w$ y $h$ para el documento «el gato come el pez».
\fuente{adaptado de Boyd y Vandenberghe~\cite{vmls}, cap.~1, ejercicio~1.11; se añade el apartado~(d).}
\end{ejercicio}

\begin{ejercicio}{2}{eje:1-normas}
Sea $x=(3,-4,12)$. Calcule $\norm{x}_1$, $\norm{x}$ y $\norm{x}_\infty$, compruebe que $\norm{x}_\infty\le\norm{x}\le\norm{x}_1$ y calcule el normalizado $\hat x$ respecto de la norma euclídea.
\end{ejercicio}

\begin{ejercicio}{3}{eje:1-cs}
Demuestre la desigualdad de Cauchy--Schwarz con el siguiente argumento, distinto del de la demostración del \cref{teo:cs}: para $v\neq 0$, la función $f(t)=\norm{u-tv}^{2}$ es un polinomio de segundo grado en $t$ que no toma valores negativos. Deduzca que $-1\le u\cdot v/(\norm{u}\norm{v})\le 1$ para $u,v\neq 0$.
\end{ejercicio}

\begin{ejercicio}{3}{eje:1-paralelogramo}
(a) Demuestre la \emph{ley del paralelogramo}: $\norm{u+v}^{2}+\norm{u-v}^{2}=2\norm{u}^{2}+2\norm{v}^{2}$ para todos $u,v\in\R^{d}$. (b) Demuestre la \emph{identidad de polarización}: $u\cdot v=\tfrac14\bigl(\norm{u+v}^{2}-\norm{u-v}^{2}\bigr)$. (c) Compruebe con $u=(1,0)$ y $v=(0,1)$ que la norma $\ell_1$ no cumple la ley del paralelogramo, y deduzca que no procede de ningún producto interior.
\end{ejercicio}

% =====================================================================
\chapter{Ángulos, ortogonalidad y similitud coseno}
\label{tema:2}
% =====================================================================

La desigualdad de Cauchy--Schwarz permite definir el ángulo entre dos vectores en cualquier dimensión. Este tema desarrolla esa geometría: ortogonalidad, proyección sobre una recta y similitud coseno, que es la medida con la que se comparan los embeddings. Se sigue la sección~1.2 de Strang~\cite{strang} y las secciones~3.3, 3.4 y~3.8 de~\cite{mml}; el comportamiento en alta dimensión procede del capítulo~2 de Blum, Hopcroft y Kannan~\cite{bhk}.

\section{El ángulo entre dos vectores}

\begin{definicion}[Ángulo]\label{def:angulo}
Sean $u,v\in\R^{d}$ no nulos. El \emph{ángulo} entre $u$ y $v$ es el único $\theta\in[0,\pi]$ tal que
\[
\cos\theta=\frac{u\cdot v}{\norm{u}\,\norm{v}}.
\]
\end{definicion}

La definición es correcta porque, por el \cref{teo:cs}, el cociente está en $[-1,1]$, y la función coseno es una biyección de $[0,\pi]$ en $[-1,1]$. La proposición siguiente muestra que, en el plano, coincide con el ángulo de la geometría elemental.

\begin{proposicion}\label{prop:ley-coseno}
Sean $u,v\in\R^{2}$ no nulos y sea $\varphi$ el ángulo geométrico que forman las flechas $u$ y $v$. Entonces $u\cdot v=\norm{u}\norm{v}\cos\varphi$.
\end{proposicion}

\begin{proof}
Los vectores $u$, $v$ y $u-v$ forman un triángulo cuyos lados miden $\norm{u}$, $\norm{v}$ y $\norm{u-v}$, y el ángulo opuesto al lado $u-v$ es $\varphi$. El teorema del coseno da $\norm{u-v}^{2}=\norm{u}^{2}+\norm{v}^{2}-2\norm{u}\norm{v}\cos\varphi$. Por otra parte, la linealidad del producto escalar da $\norm{u-v}^{2}=\norm{u}^{2}+\norm{v}^{2}-2\,u\cdot v$. Igualando ambas expresiones se obtiene el enunciado. (Si $u$ y $v$ son paralelos, el triángulo degenera y la igualdad se comprueba directamente con $u=cv$.)
\end{proof}

El signo del producto escalar indica, por tanto, si dos vectores apuntan en direcciones parecidas u opuestas (\cref{tab:signo}). El ejercicio~\ref{eje:2-tres} explora hasta dónde pueden «llevarse mal» varios vectores del plano.

\begin{table}[h]
\centering\small
\caption{Signo del producto escalar y ángulo entre dos vectores no nulos.}\label{tab:signo}
\begin{tabular}{@{}lll@{}}
\toprule
Producto escalar & Ángulo & Interpretación \\
\midrule
$u\cdot v>0$ & $0\le\theta<\pi/2$ & direcciones parecidas \\
$u\cdot v=0$ & $\theta=\pi/2$ & direcciones perpendiculares \\
$u\cdot v<0$ & $\pi/2<\theta\le\pi$ & direcciones opuestas \\
\bottomrule
\end{tabular}
\end{table}

\begin{definicion}[Ortogonalidad]
Dos vectores $u,v\in\R^{d}$ son \emph{ortogonales}, y se escribe $u\perp v$, si $u\cdot v=0$.
\end{definicion}

El vector nulo es ortogonal a todos, y es el único vector ortogonal a sí mismo, por la positividad del producto escalar. Los vectores de la base canónica son ortogonales dos a dos: $e_i\cdot e_j=0$ si $i\neq j$.

\begin{teorema}[Pitágoras]\label{teo:pitagoras}
Si $u\perp v$, entonces $\norm{u+v}^{2}=\norm{u}^{2}+\norm{v}^{2}$.
\end{teorema}

\begin{proof}
$\norm{u+v}^{2}=\norm{u}^{2}+2\,u\cdot v+\norm{v}^{2}$, y el término central es nulo.
\end{proof}

\section{Proyección sobre una recta}

\begin{proposicion}[Proyección ortogonal sobre una recta]\label{prop:proyeccion}
Sean $u,a\in\R^{d}$ con $a\neq 0$. Entre todos los vectores de la recta $\{ta: t\in\R\}$, el más próximo a $u$ es
\[
p=\frac{u\cdot a}{a\cdot a}\,a,
\]
y es el único punto de la recta tal que $u-p\perp a$. El número $(u\cdot a)/\norm{a}$ se llama \emph{proyección escalar} de $u$ sobre $a$; si $\norm{a}=1$, coincide con $u\cdot a$.
\end{proposicion}

\begin{proof}
Sea $t^{*}=(u\cdot a)/(a\cdot a)$ y $p=t^{*}a$. Entonces $(u-p)\cdot a=u\cdot a-t^{*}\,a\cdot a=0$. Recíprocamente, $(u-ta)\cdot a=0$ obliga a $t=t^{*}$. Para cualquier $t$, el vector $u-ta=(u-p)+(t^{*}-t)a$ es suma de dos vectores ortogonales, y por el \cref{teo:pitagoras}
\[
\norm{u-ta}^{2}=\norm{u-p}^{2}+(t^{*}-t)^{2}\norm{a}^{2}\ge\norm{u-p}^{2},
\]
con igualdad solo si $t=t^{*}$.
\end{proof}

\begin{figure}[h]
\centering
\begin{tikzpicture}[scale=1.15,>={Stealth[length=2.4mm]},font=\small]
  \draw[thin,gray] (-0.4,-0.2) -- (5.2,2.6);
  \draw[->,thick] (0,0) -- (2.0,1.0) node[below right] {$a$};
  \draw[->,thick] (0,0) -- (2.6,2.8) node[above left] {$u$};
  \draw[->,thick,dashed] (0,0) -- (3.2,1.6) node[below right] {$p$};
  \draw[->] (3.2,1.6) -- (2.6,2.8) node[midway,right] {$u-p$};
  \draw (3.2,1.6) ++(-0.179,-0.089) -- ++(-0.089,0.179) -- ++(0.179,0.089);
  \node at (4.9,2.15) {$\{ta\}$};
\end{tikzpicture}
\caption{Proyección ortogonal de $u$ sobre la recta generada por $a$ (\cref{prop:proyeccion}). El error $u-p$ es perpendicular a $a$.}
\label{fig:proyeccion}
\end{figure}

\begin{ejemplo}\label{ej:proyeccion}
Sean $u=(4,-1,2)$ y $a=(1,1,1)$. Como $u\cdot a=5$ y $a\cdot a=3$, la proyección es $p=\tfrac53(1,1,1)$, y el error $u-p=\tfrac13(7,-8,1)$ cumple $(u-p)\cdot a=\tfrac13(7-8+1)=0$.
\end{ejemplo}

La proyección sobre una recta es el caso más sencillo de la proyección sobre un subespacio, que conduce a los mínimos cuadrados~\cite[§4.2]{strang}; en el \cref{tema:4} se escribe como producto de matrices (ejercicio~\ref{eje:4-proyeccion}).

\section{Similitud coseno}

\begin{definicion}[Similitud coseno]\label{def:coseno}
La \emph{similitud coseno} de dos vectores no nulos $u,v\in\R^{d}$ es
\[
\cos(u,v)=\frac{u\cdot v}{\norm{u}\,\norm{v}}=\hat u\cdot\hat v\in[-1,1].
\]
\end{definicion}

\begin{proposicion}\label{prop:coseno}
Sean $u,v\in\R^{d}$ no nulos.
\begin{enumerate}
  \item $\cos(au,bv)=\cos(u,v)$ para todos $a,b>0$, y $\cos(u,-v)=-\cos(u,v)$.
  \item $\cos(u,v)=1$ si y solo si $u=cv$ con $c>0$; $\cos(u,v)=-1$ si y solo si $u=cv$ con $c<0$.
\end{enumerate}
\end{proposicion}

\begin{proof}
(i) Numerador y denominador se multiplican por $ab=\abs{ab}$; con $b=-1$, el numerador cambia de signo y el denominador no. (ii) Si $\abs{\cos(u,v)}=1$, se da la igualdad en el \cref{teo:cs}, luego $u=cv$ con $c\neq 0$, y entonces $\cos(u,v)=c\norm{v}^{2}/(\abs{c}\norm{v}^{2})$ es el signo de $c$. El recíproco es inmediato.
\end{proof}

La similitud coseno mide, por tanto, solo la dirección: $\cos((1,2),(2,4))=1$ aunque los vectores tengan longitudes distintas.

\begin{ejemplo}[Embeddings de juguete]\label{ej:juguete}
Considérense vectores de $\R^{4}$ cuyas componentes se interpretan como (realeza, masculino, persona, comida):
\begin{center}\small
\begin{tabular}{@{}lcc@{\qquad}lcc@{}}
\toprule
Palabra & Vector & Norma & Palabra & Vector & Norma \\
\midrule
rey & $(2,1,2,0)$ & $3$ & hombre & $(0,1,2,0)$ & $\sqrt5$ \\
reina & $(2,-1,2,0)$ & $3$ & mujer & $(0,-1,2,0)$ & $\sqrt5$ \\
manzana & $(0,0,0,2)$ & $2$ & & & \\
\bottomrule
\end{tabular}
\end{center}
Se obtiene $\cos(\text{rey},\text{reina})=7/9\approx 0{,}7778$ (ángulo de $38{,}94^\circ$), $\cos(\text{rey},\text{hombre})=5/(3\sqrt5)\approx 0{,}7454$, $\cos(\text{rey},\text{mujer})=3/(3\sqrt5)\approx 0{,}4472$ y $\cos(\text{rey},\text{manzana})=0$. El vecino más próximo de «rey» es «reina»: comparten realeza y humanidad y solo discrepan en una componente. En un modelo real las componentes no tienen un significado asignado de antemano, pero la comparación se hace de la misma forma.
\end{ejemplo}

\begin{proposicion}[Coseno y distancia]\label{prop:coseno-distancia}
Si $\norm{u}=\norm{v}=1$, entonces $\norm{u-v}^{2}=2-2\cos(u,v)$. En consecuencia, entre vectores normalizados, ordenar los candidatos por distancia creciente a una consulta equivale a ordenarlos por similitud coseno decreciente.
\end{proposicion}

La demostración se propone en el ejercicio~\ref{eje:2-distancia}. Sin normalizar, ambas medidas pueden discrepar: un vector muy largo en la misma dirección que la consulta tiene coseno~1 pero puede estar lejos (ejercicio~\ref{eje:2-vecinos}). Los buscadores vectoriales de los sistemas RAG (módulo~22) usan el coseno o el producto escalar sobre vectores normalizados.

\section{Geometría en alta dimensión}

En $\R^{2}$, dos direcciones elegidas al azar forman ángulos de cualquier tamaño. En dimensión alta el comportamiento es muy distinto: dos direcciones aleatorias son casi ortogonales~\cite[§2.4.2]{bhk}. El resultado siguiente lo cuantifica.

\begin{proposicion}\label{prop:alta-dim}
Sea $U$ un vector aleatorio de $\R^{d}$ con $\norm{U}=1$ cuya distribución no cambia al permutar sus componentes ni al cambiar el signo de cualquiera de ellas (por ejemplo, $U=G/\norm{G}$ con $G_1,\dots,G_d$ normales estándar independientes). Para todo vector unitario fijo $v$,
\[
\E[U\cdot v]=0,\qquad \E\bigl[(U\cdot v)^{2}\bigr]=\frac1d .
\]
En particular, la similitud coseno entre $U$ y $v$ tiene media~0 y desviación típica $1/\sqrt d$.
\end{proposicion}

\begin{proof}
Cambiar el signo de la componente $i$ no altera la distribución, luego $\E[U_i]=\E[-U_i]$ y $\E[U_i]=0$. Del mismo modo, para $i\neq j$, $\E[U_iU_j]=\E[-U_iU_j]=0$. Por la invariancia bajo permutaciones, $\E[U_i^{2}]$ es el mismo número $c$ para todo $i$, y como $\sum_iU_i^{2}=1$, se tiene $dc=1$. Entonces $\E[U\cdot v]=\sum_iv_i\E[U_i]=0$ y
\[
\E\bigl[(U\cdot v)^{2}\bigr]=\sum_{i=1}^{d}\sum_{j=1}^{d}v_iv_j\,\E[U_iU_j]=\sum_{i=1}^{d}v_i^{2}\,\frac1d=\frac1d. \qedhere
\]
\end{proof}

Si $v$ también es aleatorio e independiente de $U$, el resultado se aplica condicionando al valor de $v$. Para $d=768$, la desviación típica del coseno entre dos direcciones aleatorias es $1/\sqrt{768}\approx 0{,}0361$; para $d=12\,288$, $\approx 0{,}0090$. Una simulación con 2000 pares de vectores gaussianos da $0{,}0363$ y $0{,}0091$, respectivamente. Por eso un coseno de $0{,}3$ entre dos embeddings de dimensión 768, que equivale a más de ocho desviaciones típicas, es una señal fuerte de parecido.

\begin{observacion}
Este fenómeno permite que un espacio de dimensión $d$ contenga muchas más de $d$ direcciones casi ortogonales entre sí. La hipótesis de \emph{superposición} de la interpretabilidad (módulo~23) se apoya en esta propiedad: un modelo puede representar más conceptos que dimensiones tiene su espacio, asignando a cada uno una dirección casi ortogonal a las demás.
\end{observacion}

\begin{lecturas}
\item Strang~\cite{strang}, §1.2, problemas sobre productos escalares y ángulos.
\item Deisenroth, Faisal y Ong~\cite{mml}, §3.3 «Lengths and Distances», §3.4 «Angles and Orthogonality» y §3.8 «Orthogonal Projections».
\item Boyd y Vandenberghe~\cite{vmls}, §3.2 «Distance» y §3.4 «Angle».
\item Blum, Hopcroft y Kannan~\cite{bhk}, §2.3 «The Geometry of High Dimensions» y §2.4.2 «Volume Near the Equator».
\end{lecturas}

\section*{Ejercicios del tema 2}
\addcontentsline{toc}{section}{Ejercicios del tema 2}

\begin{ejercicio}{1}{eje:2-escala}
Calcule el producto escalar y la similitud coseno entre $(1,2)$ y $(2,4)$, y entre $(1,2)$ y $(10,20)$. ¿Qué indica la comparación sobre lo que mide cada cantidad?
\end{ejercicio}

\begin{ejercicio}{1}{eje:2-vf}
Justifique si son verdaderas o falsas: (a) si $u\cdot v=0$, uno de los dos vectores es nulo; (b) $\cos(u,u)=1$ para todo $u\neq 0$; (c) si $\cos(u,v)=1$, entonces $u=v$.
\end{ejercicio}

\begin{ejercicio}{2}{eje:2-angulo}
Sean $u=(2,-1,2)$ y $v=(1,2,2)$. Calcule $u\cdot v$, $\norm{u}$, $\norm{v}$, $\cos(u,v)$ y el ángulo entre ambos en grados.
\end{ejercicio}

\begin{ejercicio}{2}{eje:2-vecinos}
Sean $a=(1,1)$, $b=(3,3)$ y $c=(1,-1)$. Calcule $\norm{a-b}$, $\norm{a-c}$, $\cos(a,b)$ y $\cos(a,c)$. ¿Qué vecino de $a$ elige cada medida? ¿Qué ocurre si antes se normalizan los tres vectores?
\end{ejercicio}

\begin{ejercicio}{2}{eje:2-normalizar}
Normalice $x=(3,4)$. Con $p=\hat x$ y $r=(1,0)$, compruebe numéricamente que $\norm{p-r}^{2}=2-2\cos(p,r)$.
\end{ejercicio}

\begin{ejercicio}{2}{eje:2-tres}
(a) ¿Pueden tres vectores $u,v,w\in\R^{2}$ cumplir a la vez $u\cdot v<0$, $v\cdot w<0$ y $u\cdot w<0$? Si es posible, dé un ejemplo con vectores unitarios y calcule el ángulo entre cada par. (b) ¿Pueden hacerlo cuatro vectores de $\R^{2}$, con los seis productos escalares negativos?
\fuente{adaptado de Strang~\cite{strang}, §1.2, problema~28, propuesto en el curso 18.06 del MIT~\cite{mit1806}, primavera de 2010, hoja de problemas~1.}
\end{ejercicio}

\begin{ejercicio}{3}{eje:2-distancia}
Demuestre la \cref{prop:coseno-distancia}: si $\norm{u}=\norm{v}=1$, entonces $\norm{u-v}^{2}=2-2\cos(u,v)$. Deduzca que, si $q,a,b$ son unitarios y $\cos(q,a)>\cos(q,b)$, entonces $\norm{q-a}<\norm{q-b}$. ¿Qué distancia corresponde a $\cos(u,v)=-1$?
\end{ejercicio}

% =====================================================================
\chapter{Matrices y aplicaciones lineales}
\label{tema:3}
% =====================================================================

Una matriz es una tabla de números, pero su papel en el álgebra lineal es el de una función: transforma vectores en vectores respetando las combinaciones lineales. Este tema establece esa correspondencia entre matrices y aplicaciones lineales. Se sigue la sección~2.1 de Strang~\cite{strang} y la sesión «The Geometry of Linear Equations» de 18.06SC~\cite{mit1806sc}; las definiciones proceden de~\cite[§2.2 y §2.7]{mml} y de~\cite[cap.~6 y §8.1]{vmls}.

\section{Matrices}

\begin{definicion}[Matriz]
Sean $m,n\ge 1$. Una \emph{matriz} real de tamaño $m\times n$ es una tabla $A=(A_{ij})$ de $mn$ números reales dispuestos en $m$ filas y $n$ columnas; $A_{ij}$ es la entrada de la fila $i$ y la columna $j$. El conjunto de estas matrices se denota $\R^{m\times n}$. Se escribe $a_j\in\R^{m}$ para la columna $j$ y $A=[\,a_1\ \cdots\ a_n\,]$.
\end{definicion}

Un vector de $\R^{d}$ se identifica con una matriz $d\times 1$ (vector columna). Las matrices del mismo tamaño se suman y se multiplican por escalares entrada a entrada, y con esas operaciones $\R^{m\times n}$ es un espacio vectorial.

\begin{definicion}[Traspuesta e identidad]
La \emph{traspuesta} de $A\in\R^{m\times n}$ es la matriz $A^{\top}\in\R^{n\times m}$ con $(A^{\top})_{ij}=A_{ji}$: las filas de $A$ son las columnas de $A^{\top}$. Una matriz cuadrada es \emph{simétrica} si $A^{\top}=A$. La \emph{identidad} $I_n\in\R^{n\times n}$ tiene unos en la diagonal y ceros fuera de ella; sus columnas son los vectores $e_1,\dots,e_n$ de la base canónica.
\end{definicion}

Si $u\in\R^{d}$ es una columna, $u^{\top}$ es una matriz $1\times d$ (vector fila). Esta notación permite escribir el producto escalar como $u\cdot v=u^{\top}v$.

\section{El producto de una matriz por un vector}

\begin{definicion}[Producto matriz-vector]\label{def:matvec}
Sean $A\in\R^{m\times n}$ y $x\in\R^{n}$. El producto $Ax\in\R^{m}$ es el vector de componentes
\[
(Ax)_i=\sum_{j=1}^{n}A_{ij}x_j,\qquad i=1,\dots,m.
\]
El producto solo está definido si el número de columnas de $A$ coincide con el número de componentes de $x$.
\end{definicion}

\begin{proposicion}[Dos lecturas del producto]\label{prop:filas-columnas}
Sea $A\in\R^{m\times n}$ con filas $r_1^{\top},\dots,r_m^{\top}$ y columnas $a_1,\dots,a_n$. Para todo $x\in\R^{n}$:
\begin{enumerate}
  \item (por filas) la componente $i$ de $Ax$ es el producto escalar $r_i\cdot x$;
  \item (por columnas) $Ax=x_1a_1+x_2a_2+\dots+x_na_n$ es una combinación lineal de las columnas de $A$.
\end{enumerate}
\end{proposicion}

\begin{proof}
(i) La fila $i$ de $A$ es $(A_{i1},\dots,A_{in})$, y su producto escalar con $x$ es $\sum_jA_{ij}x_j$. (ii) La componente $i$ de $\sum_jx_ja_j$ es $\sum_jx_jA_{ij}$, que es la misma suma.
\end{proof}

\begin{ejemplo}\label{ej:matvec}
Sean $W=\begin{pmatrix}1&0&2\\0&-1&1\end{pmatrix}$ y $x=(3,1,2)$. Por filas, $(Wx)_1=3+0+4=7$ y $(Wx)_2=0-1+2=1$. Por columnas, $3(1,0)+1(0,-1)+2(2,1)=(7,1)$. La matriz transforma un vector de $\R^{3}$ en uno de $\R^{2}$.
\end{ejemplo}

Las dos lecturas corresponden a las dos imágenes de un sistema de ecuaciones lineales $Ax=b$ que Strang llama imagen por filas e imagen por columnas~\cite[§2.1]{strang}: por filas, cada ecuación es un hiperplano y la solución es su intersección; por columnas, se buscan los coeficientes con los que las columnas de $A$ se combinan para dar $b$.

\section{Aplicaciones lineales}

\begin{definicion}[Aplicación lineal]\label{def:lineal}
Una aplicación $T\colon\R^{n}\to\R^{m}$ es \emph{lineal} si
\[
T(au+bv)=a\,T(u)+b\,T(v)\qquad\text{para todos } u,v\in\R^{n},\ a,b\in\R.
\]
\end{definicion}

Tomando $a=b=0$ se obtiene $T(0)=0$: una aplicación lineal fija el origen. Por inducción, $T$ transforma cualquier combinación lineal en la combinación lineal de las imágenes con los mismos coeficientes.

\begin{proposicion}\label{prop:matriz-lineal}
Para toda $A\in\R^{m\times n}$, la aplicación $x\mapsto Ax$ es lineal.
\end{proposicion}

\begin{proof}
La componente $i$ de $A(au+bv)$ es $\sum_jA_{ij}(au_j+bv_j)=a\sum_jA_{ij}u_j+b\sum_jA_{ij}v_j$, que es la componente $i$ de $aAu+bAv$.
\end{proof}

El resultado recíproco es el que justifica la frase «una matriz es una aplicación lineal».

\begin{teorema}[Matriz de una aplicación lineal]\label{teo:representacion}
Sea $T\colon\R^{n}\to\R^{m}$ lineal y sea $A\in\R^{m\times n}$ la matriz cuyas columnas son $T(e_1),\dots,T(e_n)$. Entonces $T(x)=Ax$ para todo $x\in\R^{n}$, y $A$ es la única matriz con esta propiedad.
\end{teorema}

\begin{proof}
Todo $x\in\R^{n}$ se escribe $x=x_1e_1+\dots+x_ne_n$. Por la linealidad de $T$ y la \cref{prop:filas-columnas}(ii),
\[
T(x)=x_1T(e_1)+\dots+x_nT(e_n)=x_1a_1+\dots+x_na_n=Ax.
\]
Si $B$ también cumple $T(x)=Bx$ para todo $x$, entonces la columna $j$ de $B$ es $Be_j=T(e_j)=a_j$, luego $B=A$.
\end{proof}

El \cref{teo:representacion} da una receta práctica: para conocer una aplicación lineal basta saber adónde envía los vectores de la base canónica, y esas imágenes son las columnas de su matriz. El \cref{tab:transformaciones} recoge las transformaciones del plano más habituales, y la \cref{fig:cizalla} muestra el efecto de una cizalla sobre el cuadrado unidad.

\begin{table}[h]
\centering\small
\caption{Transformaciones lineales del plano. Las columnas de cada matriz son las imágenes de $e_1$ y $e_2$.}\label{tab:transformaciones}
\begin{tabular}{@{}lccc@{}}
\toprule
Transformación & Matriz & $e_1\mapsto$ & $e_2\mapsto$ \\
\midrule
Rotación de ángulo $\theta$ & $R_\theta=\begin{pmatrix}\cos\theta&-\sin\theta\\ \sin\theta&\cos\theta\end{pmatrix}$ & $(\cos\theta,\sin\theta)$ & $(-\sin\theta,\cos\theta)$ \\[10pt]
Rotación de $90^\circ$ & $\begin{pmatrix}0&-1\\1&0\end{pmatrix}$ & $(0,1)$ & $(-1,0)$ \\[8pt]
Escalado & $\begin{pmatrix}2&0\\0&0{,}5\end{pmatrix}$ & $(2,0)$ & $(0;\,0{,}5)$ \\[8pt]
Cizalla & $\begin{pmatrix}1&1\\0&1\end{pmatrix}$ & $(1,0)$ & $(1,1)$ \\[8pt]
Simetría respecto de $y=x$ & $\begin{pmatrix}0&1\\1&0\end{pmatrix}$ & $(0,1)$ & $(1,0)$ \\[8pt]
Proyección sobre el eje $x$ & $\begin{pmatrix}1&0\\0&0\end{pmatrix}$ & $(1,0)$ & $(0,0)$ \\
\bottomrule
\end{tabular}
\end{table}

\begin{figure}[h]
\centering
\begin{tikzpicture}[scale=1.5,>={Stealth[length=2.2mm]},font=\small]
  \draw[->,gray] (-0.3,0) -- (2.4,0);
  \draw[->,gray] (0,-0.3) -- (0,1.4);
  \fill[gray!15] (0,0) -- (1,0) -- (2,1) -- (1,1) -- cycle;
  \draw[thick] (0,0) -- (1,0) -- (2,1) -- (1,1) -- cycle;
  \draw[dashed] (0,0) rectangle (1,1);
  \draw[->,very thick] (0,0) -- (1,0) node[below] {$e_1=Ae_1$};
  \draw[->,very thick] (0,0) -- (0,1) node[left] {$e_2$};
  \draw[->,very thick] (0,0) -- (1,1) node[above] {$Ae_2$};
\end{tikzpicture}
\caption{La cizalla $A=\begin{pmatrix}1&1\\0&1\end{pmatrix}$ transforma el cuadrado unidad (discontinuo) en un paralelogramo. La imagen de $e_1$ no cambia y la de $e_2$ es la segunda columna de $A$.}
\label{fig:cizalla}
\end{figure}

\begin{proposicion}[Las rotaciones conservan la norma]\label{prop:rotacion}
Para todo $\theta\in\R$ y todo $x\in\R^{2}$, $\norm{R_\theta x}=\norm{x}$.
\end{proposicion}

\begin{proof}
Con $x=(x_1,x_2)$, $R_\theta x=(x_1\cos\theta-x_2\sin\theta,\ x_1\sin\theta+x_2\cos\theta)$. Desarrollando los cuadrados, los términos cruzados $\mp 2x_1x_2\sin\theta\cos\theta$ se cancelan y queda
$\norm{R_\theta x}^{2}=(x_1^{2}+x_2^{2})(\cos^{2}\theta+\sin^{2}\theta)=\norm{x}^{2}$.
\end{proof}

Las rotaciones conservan también los productos escalares y, por tanto, los ángulos (ejercicio~\ref{eje:4-ortogonal}). Los escalados no: con la matriz de escalado del \cref{tab:transformaciones}, $e_1+e_2$, de norma $\sqrt2$, pasa a $(2;\,0{,}5)$, de norma $\sqrt{4{,}25}$.

\begin{observacion}[Matrices en los modelos de lenguaje]\label{obs:qkv}
Casi todo el cálculo de un Transformer consiste en aplicar matrices aprendidas a vectores. En la atención (módulo~13), cada embedding $x\in\R^{d}$ se transforma en una \emph{consulta} $q=W_Qx$, una \emph{clave} $k=W_Kx$ y un \emph{valor} $v=W_Vx$, con $W_Q,W_K,W_V$ matrices $d_k\times d$. Estas matrices seleccionan qué rasgos del embedding intervienen en cada comparación.
\end{observacion}

\begin{observacion}[Matrices de transición]
La matriz de transición $P$ de un modelo de bigramas (módulo~1, tema~2) es estocástica por filas. Si $\pi_t\in\R^{K}$ es la distribución del token en la posición $t$, la de la posición siguiente es $\pi_{t+1}=P^{\top}\pi_t$, cuya componente $j$ es $\sum_iP_{ij}(\pi_t)_i$. La matriz $P^{\top}$ tiene columnas que suman~1; es la convención que usa Strang para las \emph{matrices de Markov} (ejercicio~\ref{eje:3-markov}).
\end{observacion}

\begin{lecturas}
\item Strang~\cite{strang}, §2.1, imágenes por filas y por columnas de $Ax=b$.
\item MIT 18.06SC~\cite{mit1806sc}, unidad~I, sesión «The Geometry of Linear Equations».
\item Deisenroth, Faisal y Ong~\cite{mml}, §2.2 «Matrices» y §2.7 «Linear Mappings».
\item Boyd y Vandenberghe~\cite{vmls}, §6.4 «Matrix-vector multiplication» y §8.1 «Linear and affine functions».
\end{lecturas}

\section*{Ejercicios del tema 3}
\addcontentsline{toc}{section}{Ejercicios del tema 3}

\begin{ejercicio}{1}{eje:3-lineal}
(a) Indique cuáles de las siguientes aplicaciones de $\R^{2}$ en $\R^{2}$ son lineales y, para las que lo sean, escriba su matriz: $f(x)=(x_1+x_2,\,3x_1)$, $g(x)=(x_1+1,\,x_2)$, $h(x)=(x_1x_2,\,x_2)$. (b) Una matriz $W$ es de tamaño $5\times 3$. ¿De qué espacio a qué espacio transforma vectores?
\end{ejercicio}

\begin{ejercicio}{2}{eje:3-dos-lecturas}
Sean $M=\begin{pmatrix}1&0&2\\2&1&0\\0&3&1\end{pmatrix}$ y $x=(2,-1,3)$. Calcule $Mx$ primero por filas y después como combinación lineal de las columnas. ¿Qué relación geométrica hay entre la tercera fila de $M$ y $x$?
\fuente{adaptado de MIT 18.06SC~\cite{mit1806sc}, unidad~I, sesión «The Geometry of Linear Equations», problema~1.2 (matriz y vector cambiados).}
\end{ejercicio}

\begin{ejercicio}{2}{eje:3-base}
Una aplicación lineal $T\colon\R^{2}\to\R^{2}$ cumple $T(e_1)=(1,2)$ y $T(e_2)=(-1,1)$. Escriba su matriz y calcule $T(3,-2)$ de dos formas: con la matriz y con la linealidad.
\end{ejercicio}

\begin{ejercicio}{2}{eje:3-rotacion}
Escriba la matriz $R_{60^\circ}$ y calcule $R_{60^\circ}(2,0)$ y $R_{60^\circ}(1,1)$. Compruebe que las normas se conservan.
\end{ejercicio}

\begin{ejercicio}{2}{eje:3-markov}
Sea $A=\begin{pmatrix}0{,}8&0{,}3\\0{,}2&0{,}7\end{pmatrix}$, cuyas columnas suman~1, y sea $u_0=(1,0)$. (a) Calcule $u_1=Au_0$, $u_2=Au_1$ y $u_3=Au_2$. ¿Qué propiedad comparten las componentes de todos ellos? (b) Compruebe que $s=(0{,}6;\,0{,}4)$ cumple $As=s$. (c) Demuestre que, si las columnas de una matriz $A\in\R^{n\times n}$ suman~1, entonces $\unos^{\top}(Au)=\unos^{\top}u$ para todo $u\in\R^{n}$.
\fuente{adaptado de Strang~\cite{strang}, §2.1, problemas~29 y~30, propuestos en el curso 18.06 del MIT~\cite{mit1806}, primavera de 2010, hoja de problemas~1; se añade el apartado~(c).}
\end{ejercicio}

\begin{ejercicio}{3}{eje:3-rectas}
Sea $T\colon\R^{n}\to\R^{m}$ lineal. (a) Demuestre que $T$ transforma la recta $\{p+td: t\in\R\}$ en una recta o en un punto, y caracterice cuándo ocurre lo segundo. (b) Demuestre que $T$ conserva los puntos medios: $T\bigl(\tfrac12(u+v)\bigr)=\tfrac12\bigl(T(u)+T(v)\bigr)$. (c) Sea $b\neq 0$. Demuestre que la traslación $x\mapsto x+b$ cumple (a) y (b) pero no es lineal.
\end{ejercicio}

% =====================================================================
\chapter{El producto de matrices}
\label{tema:4}
% =====================================================================

El producto de matrices se define de modo que represente la composición de aplicaciones lineales. De esa definición se siguen sus propiedades (asociativa, distributivas, comportamiento con la traspuesta) y su principal diferencia con el producto de números: no es conmutativo. El tema sigue la sección~2.4 de Strang~\cite{strang} y la sesión «Multiplication and Inverse Matrices» de 18.06SC~\cite{mit1806sc}; la composición y el coste computacional se tratan como en~\cite[§10.1 y §10.2]{vmls}.

\section{Definición y composición}

\begin{definicion}[Producto de matrices]\label{def:producto}
Sean $A\in\R^{m\times n}$ y $B\in\R^{n\times p}$. El producto $AB\in\R^{m\times p}$ es la matriz de entradas
\[
(AB)_{ij}=\sum_{k=1}^{n}A_{ik}B_{kj},\qquad 1\le i\le m,\ 1\le j\le p.
\]
La entrada $(i,j)$ es el producto escalar de la fila $i$ de $A$ por la columna $j$ de $B$. El producto solo está definido si el número de columnas de $A$ coincide con el número de filas de $B$: $(m\times n)(n\times p)=m\times p$.
\end{definicion}

\begin{teorema}[El producto es la composición]\label{teo:composicion}
Sean $A\in\R^{m\times n}$ y $B\in\R^{n\times p}$. Para todo $x\in\R^{p}$, $(AB)x=A(Bx)$. Es decir, la matriz de la aplicación «aplicar primero $B$ y después $A$» es $AB$.
\end{teorema}

\begin{proof}
La componente $i$ de $A(Bx)$ es
\[
\sum_{k=1}^{n}A_{ik}(Bx)_k=\sum_{k=1}^{n}A_{ik}\sum_{j=1}^{p}B_{kj}x_j=\sum_{j=1}^{p}\Bigl(\sum_{k=1}^{n}A_{ik}B_{kj}\Bigr)x_j=\sum_{j=1}^{p}(AB)_{ij}x_j,
\]
que es la componente $i$ de $(AB)x$. Se ha usado solo que las sumas finitas pueden reordenarse.
\end{proof}

El orden de escritura es el de la composición de funciones, $(A\circ B)(x)=A(B(x))$: en $AB$ actúa primero la matriz de la derecha.

\begin{ejemplo}\label{ej:intercambio}
Sean $A=\begin{pmatrix}1&2\\3&4\end{pmatrix}$ y $B=\begin{pmatrix}0&1\\1&0\end{pmatrix}$. Entonces
\[
AB=\begin{pmatrix}2&1\\4&3\end{pmatrix},\qquad BA=\begin{pmatrix}3&4\\1&2\end{pmatrix}.
\]
Multiplicar por $B$ a la derecha intercambia las columnas de $A$; multiplicar por $B$ a la izquierda intercambia sus filas.
\end{ejemplo}

\section{Cuatro formas de calcular un producto}

La \cref{def:producto} calcula el producto entrada a entrada. Strang subraya otras tres formas equivalentes~\cite[§2.4]{strang}, cada una útil en contextos distintos.

\begin{proposicion}\label{prop:cuatro}
Sean $A\in\R^{m\times n}$ con columnas $a_1,\dots,a_n$ y filas $r_1^{\top},\dots,r_m^{\top}$, y $B\in\R^{n\times p}$ con columnas $b_1,\dots,b_p$ y filas $s_1^{\top},\dots,s_n^{\top}$. Entonces:
\begin{enumerate}
  \item (por columnas) la columna $j$ de $AB$ es $Ab_j$;
  \item (por filas) la fila $i$ de $AB$ es $r_i^{\top}B$;
  \item (columnas por filas) $AB=\sum_{k=1}^{n}a_k\,s_k^{\top}$, suma de $n$ matrices $m\times p$.
\end{enumerate}
\end{proposicion}

\begin{proof}
(i) La componente $i$ de $Ab_j$ es $\sum_kA_{ik}B_{kj}=(AB)_{ij}$. (ii) La componente $j$ de $r_i^{\top}B$ es $\sum_kA_{ik}B_{kj}$. (iii) La entrada $(i,j)$ de la matriz $a_ks_k^{\top}$ es $A_{ik}B_{kj}$; sumando en $k$ se obtiene $(AB)_{ij}$.
\end{proof}

Una matriz de la forma $a\,s^{\top}$, con $a\in\R^{m}$ y $s\in\R^{p}$, se llama \emph{producto exterior}; todas sus filas son múltiplos de $s^{\top}$. La forma (iii) descompone un producto de matrices en suma de productos exteriores y es la base de la descomposición en valores singulares y de las aproximaciones de bajo rango del módulo~3. En el \cref{ej:intercambio}, $a_1s_1^{\top}+a_2s_2^{\top}=\begin{pmatrix}0&1\\0&3\end{pmatrix}+\begin{pmatrix}2&0\\4&0\end{pmatrix}=AB$.

\section{Propiedades algebraicas}

\begin{proposicion}\label{prop:algebra}
Siempre que los tamaños permitan formar los productos:
\begin{enumerate}
  \item (asociativa) $(AB)C=A(BC)$;
  \item (distributivas) $A(B+C)=AB+AC$ y $(B+C)A=BA+CA$;
  \item (identidad) $I_mA=A=AI_n$ si $A\in\R^{m\times n}$;
  \item (traspuesta) $(AB)^{\top}=B^{\top}A^{\top}$.
\end{enumerate}
\end{proposicion}

\begin{proof}
(i) Por el \cref{teo:composicion}, para todo $x$, $((AB)C)x=(AB)(Cx)=A(B(Cx))$ y $(A(BC))x=A((BC)x)=A(B(Cx))$. Dos matrices que actúan igual sobre todo vector son iguales (basta tomar $x=e_j$ para ver que tienen las mismas columnas). (ii) $\sum_kA_{ik}(B_{kj}+C_{kj})=\sum_kA_{ik}B_{kj}+\sum_kA_{ik}C_{kj}$, y análogamente por la derecha. (iii) Solo $k=i$ contribuye en $\sum_k(I_m)_{ik}A_{kj}$. (iv) $((AB)^{\top})_{ij}=(AB)_{ji}=\sum_kA_{jk}B_{ki}=\sum_k(B^{\top})_{ik}(A^{\top})_{kj}=(B^{\top}A^{\top})_{ij}$.
\end{proof}

La distributiva por la izquierda y la distributiva por la derecha son propiedades distintas: no puede pasarse de una a otra cambiando el lado (ejercicio~\ref{eje:4-distributiva}), porque el producto no es conmutativo.

\begin{observacion}[No conmutatividad]\label{obs:no-conmuta}
En general $AB\neq BA$, incluso si ambos productos están definidos y son del mismo tamaño. Sean $R=\begin{pmatrix}0&-1\\1&0\end{pmatrix}$ (rotación de $90^\circ$), $S=\begin{pmatrix}2&0\\0&0{,}5\end{pmatrix}$ (escalado) y $x=(3,1)$. Escalar y después rotar da $R(Sx)=R(6;\,0{,}5)=(-0{,}5;\,6)$; rotar y después escalar da $S(Rx)=S(-1,3)=(-2;\,1{,}5)$. El orden de las transformaciones importa.
\end{observacion}

\begin{definicion}[Matriz ortogonal]\label{def:ortogonal}
Una matriz $Q\in\R^{n\times n}$ es \emph{ortogonal} si $Q^{\top}Q=I_n$, es decir, si sus columnas son unitarias y ortogonales dos a dos.
\end{definicion}

\begin{proposicion}\label{prop:ortogonal}
Si $Q$ es ortogonal, entonces $(Qu)\cdot(Qv)=u\cdot v$ para todos $u,v\in\R^{n}$. En consecuencia, $Q$ conserva normas, distancias y ángulos.
\end{proposicion}

\begin{proof}
Por la \cref{prop:algebra}(iv), $(Qu)^{\top}(Qv)=u^{\top}Q^{\top}Qv=u^{\top}I_nv=u^{\top}v$. Las normas, las distancias y los ángulos se expresan mediante productos escalares.
\end{proof}

Las rotaciones $R_\theta$ son ortogonales (ejercicio~\ref{eje:4-ortogonal}), lo que generaliza la \cref{prop:rotacion}. Las codificaciones posicionales rotatorias (RoPE) del módulo~14 aplican precisamente rotaciones de este tipo a las consultas y las claves.

\section{Coste y matriz de Gram}

\begin{proposicion}[Coste del producto]\label{prop:coste}
Calcular $AB$ con $A\in\R^{m\times n}$ y $B\in\R^{n\times p}$ mediante la \cref{def:producto} requiere $mnp$ multiplicaciones y $mp(n-1)$ sumas.
\end{proposicion}

\begin{proof}
Cada una de las $mp$ entradas es un producto escalar en $\R^{n}$, que requiere $n$ multiplicaciones y $n-1$ sumas.
\end{proof}

Aunque $(AB)C=A(BC)$, el coste de los dos órdenes puede ser muy distinto (ejercicio~\ref{eje:4-coste}). Elegir bien el orden de los paréntesis es una optimización habitual en el cálculo con tensores.

\begin{definicion}[Matriz de Gram]\label{def:gram}
Sea $X\in\R^{n\times d}$ la matriz cuyas filas son los vectores $x_1^{\top},\dots,x_n^{\top}$. Su \emph{matriz de Gram} es $G=XX^{\top}\in\R^{n\times n}$, con $G_{ij}=x_i\cdot x_j$.
\end{definicion}

La matriz de Gram es simétrica, porque $G^{\top}=(XX^{\top})^{\top}=XX^{\top}$, y su diagonal contiene las normas al cuadrado, $G_{ii}=\norm{x_i}^{2}$. Un único producto de matrices calcula los $n^{2}$ productos escalares entre $n$ vectores.

\begin{ejemplo}\label{ej:gram}
Con los embeddings del \cref{ej:juguete}, sea $X$ la matriz $3\times 4$ de filas «rey», «reina» y «manzana». Entonces
\[
XX^{\top}=\begin{pmatrix}9&7&0\\7&9&0\\0&0&4\end{pmatrix}.
\]
La entrada $(1,2)$ es $\text{rey}\cdot\text{reina}=7$, y dividiendo por $\sqrt{G_{11}G_{22}}=9$ se recupera la similitud coseno $7/9$.
\end{ejemplo}

\begin{observacion}
En la atención, las consultas y las claves de los $n$ tokens de una secuencia se apilan como filas de matrices $Q,K\in\R^{n\times d_k}$, y todas las puntuaciones se calculan con un único producto $QK^{\top}\in\R^{n\times n}$, que generaliza la matriz de Gram. Las GPU están diseñadas para ejecutar este tipo de productos con gran rapidez (módulo~17).
\end{observacion}

\begin{lecturas}
\item Strang~\cite{strang}, §2.4, reglas de las operaciones con matrices, y §2.7, traspuestas.
\item MIT 18.06SC~\cite{mit1806sc}, unidad~I, sesión «Multiplication and Inverse Matrices».
\item Deisenroth, Faisal y Ong~\cite{mml}, §2.2 «Matrices» y §3.9 «Rotations».
\item Boyd y Vandenberghe~\cite{vmls}, §10.1 «Matrix-matrix multiplication» y §10.2 «Composition of linear functions».
\end{lecturas}

\section*{Ejercicios del tema 4}
\addcontentsline{toc}{section}{Ejercicios del tema 4}

\begin{ejercicio}{1}{eje:4-formas}
(a) ¿Qué tamaño tiene el producto de una matriz $3\times 4$ por otra $4\times 2$? ¿Puede multiplicarse una matriz $2\times 3$ por otra $2\times 3$? (b) Para $u,v\in\R^{3}$, ¿qué tamaños tienen $u^{\top}v$ y $uv^{\top}$? (c) Justifique si es verdadero o falso que $AB=BA$ para todo par de matrices cuadradas del mismo orden.
\end{ejercicio}

\begin{ejercicio}{2}{eje:4-rs}
Con $R$ y $S$ de la \cref{obs:no-conmuta}, calcule $RS$ y $SR$, aplíquelas a $x=(3,1)$ y compruebe que se obtiene $R(Sx)$ y $S(Rx)$, respectivamente.
\end{ejercicio}

\begin{ejercicio}{2}{eje:4-ab-ba}
Sean $A=\begin{pmatrix}1&0&2\\0&1&-1\end{pmatrix}$ y $B=\begin{pmatrix}1&2\\0&1\\3&0\end{pmatrix}$. Calcule $AB$ y $BA$ e indique sus tamaños. ¿Puede calcularse $AA$?
\end{ejercicio}

\begin{ejercicio}{2}{eje:4-distributiva}
Sean $A=\begin{pmatrix}2&1\\1&3\end{pmatrix}$, $B=\begin{pmatrix}1&0\\0&0\end{pmatrix}$ y $C=\begin{pmatrix}0&0\\2&1\end{pmatrix}$. Calcule $AB+AC$ y $A(B+C)$ y compárelas. Calcule también $BA+CA$. ¿Coincide con $A(B+C)$?
\fuente{adaptado de MIT 18.06SC~\cite{mit1806sc}, unidad~I, sesión «Multiplication and Inverse Matrices», problema~3.1 (matrices cambiadas).}
\end{ejercicio}

\begin{ejercicio}{2}{eje:4-coste}
Sean $A\in\R^{3\times 5}$, $B\in\R^{5\times 2}$ y $C\in\R^{2\times 4}$. (a) Con la \cref{prop:coste}, calcule el número de multiplicaciones de $(AB)C$ y de $A(BC)$. (b) Para $u,v,w\in\R^{N}$, compare el coste de $(u^{\top}v)\,w^{\top}$ y de $u^{\top}(vw^{\top})$, y evalúelo para $N=4096$.
\fuente{adaptado de Strang~\cite{strang}, §2.4, problema~36, propuesto en el curso 18.06 del MIT~\cite{mit1806}, primavera de 2010, hoja de problemas~1 (dimensiones cambiadas).}
\end{ejercicio}

\begin{ejercicio}{2}{eje:4-proyeccion}
Sea $A\in\R^{3\times 2}$ la matriz formada por las dos primeras columnas de $I_3$, y sea $b=(4,-1,2)$. (a) Como las columnas de $A$ son unitarias y ortogonales, la matriz de proyección sobre el plano que generan es $P=AA^{\top}$. Calcule $P$ y $Pb$. (b) Compruebe que $P^{2}=P$ y que $(I-P)^{2}=I-P$. ¿Sobre qué proyecta $I-P$? (c) Escriba la proyección sobre la recta de $a=(1,1,1)$ de la \cref{prop:proyeccion} como $P_a\,b$ con $P_a=aa^{\top}/(a^{\top}a)$ y compruebe el resultado del \cref{ej:proyeccion}.
\fuente{adaptado de MIT 18.06SC~\cite{mit1806sc}, unidad~II, sesión «Projections onto Subspaces», problemas~15.1 y~15.2 (dimensión y vector cambiados).}
\end{ejercicio}

\begin{ejercicio}{3}{eje:4-ortogonal}
(a) Demuestre que $R_\theta^{\top}R_\theta=I_2$ para todo $\theta$. (b) Demuestre que $R_\theta R_\varphi=R_{\theta+\varphi}$ y deduzca que dos rotaciones del plano siempre conmutan. (c) Demuestre que el producto de dos matrices ortogonales del mismo orden es ortogonal. (d) Contraste (b) con la \cref{obs:no-conmuta}: ¿qué propiedad de $S$ impide que conmute con $R$?
\end{ejercicio}

% =====================================================================
\chapter{Embeddings, analogías y atención}
\label{tema:5}
% =====================================================================

Este tema aplica las herramientas anteriores a la entrada y la salida de un modelo de lenguaje. La capa de embeddings es una matriz; leer el vector de un token es multiplicar por un vector \emph{one-hot}; las relaciones semánticas aparecen, de forma aproximada, como direcciones del espacio; y la atención compara tokens mediante productos escalares. La presentación sigue el capítulo~5 de Jurafsky y Martin~\cite{jm3} y la clase~2 de CS224N~\cite{cs224n}, junto con los artículos de Mikolov y coautores~\cite{mikolov-a,mikolov-b} y de Vaswani y coautores~\cite{vaswani}.

\section{La matriz de embeddings}

En este tema se usa la notación del módulo~1: $\V$ es el vocabulario, $K=\abs{\V}$ su tamaño y los tokens se identifican con sus índices $0,1,\dots,K-1$. Las filas de las matrices indexadas por tokens se numeran también desde~0.

\begin{definicion}[Matriz de embeddings]\label{def:embedding}
Sea $d\ge 1$. Una \emph{matriz de embeddings} es una matriz $E\in\R^{K\times d}$. El \emph{embedding} del token $i$ es la fila $i$ de $E$, que se denota $E_{i,:}\in\R^{d}$ (escrita como columna cuando se opera con ella).
\end{definicion}

\begin{definicion}[Vector one-hot]
El vector \emph{one-hot} del token $i$ es el vector $e_i\in\R^{K}$ que tiene un~1 en la posición $i$ y ceros en el resto.
\end{definicion}

\begin{proposicion}\label{prop:onehot}
Para todo token $i$, $e_i^{\top}E=E_{i,:}^{\top}$; es decir, multiplicar el vector one-hot de un token por la matriz de embeddings selecciona su fila. Si $x_{1:n}$ es una secuencia de tokens y $O\in\R^{n\times K}$ es la matriz cuyas filas son $e_{x_1}^{\top},\dots,e_{x_n}^{\top}$, entonces $X=OE\in\R^{n\times d}$ es la matriz cuya fila $t$ es el embedding de $x_t$.
\end{proposicion}

\begin{proof}
La componente $j$ de $e_i^{\top}E$ es $\sum_k(e_i)_kE_{kj}$. Todos los sumandos con $k\neq i$ se anulan y queda $E_{ij}$. La segunda afirmación se sigue de la \cref{prop:cuatro}(ii): la fila $t$ de $OE$ es $e_{x_t}^{\top}E$.
\end{proof}

\begin{ejemplo}\label{ej:embedding}
Sean $K=4$, $d=3$ y
\[
E=\begin{pmatrix}0{,}1&0{,}2&0\\1&0&-1\\0{,}5&0{,}5&0{,}5\\-0{,}2&0{,}3&0{,}1\end{pmatrix}.
\]
El token~2 tiene vector one-hot $e_2=(0,0,1,0)$ y $e_2^{\top}E=(0{,}5;\,0{,}5;\,0{,}5)$. La secuencia $(2,0,2)$ produce una matriz $X\in\R^{3\times 3}$ cuyas filas primera y tercera coinciden: sin información de posición, el modelo no distingue las dos apariciones del token~2 (módulo~14).
\end{ejemplo}

\begin{observacion}[Indexar en lugar de multiplicar]
El producto $e_i^{\top}E$ requiere $Kd$ multiplicaciones, casi todas por cero; leer directamente la fila $i$ solo copia $d$ números. Las bibliotecas de redes neuronales implementan la capa de embeddings como una indexación, aunque matemáticamente sea el producto de la \cref{prop:onehot}. En GPT-2, con $K=50\,257$~\cite[§2.3]{radford} y $d=768$, la matriz $E$ tiene $50\,257\cdot 768=38\,597\,376$ parámetros, unos $154$~MB en precisión de 32 bits.
\end{observacion}

\section{De dónde salen los embeddings}

Los embeddings no se diseñan a mano como en el \cref{ej:juguete}: se inicializan al azar y se ajustan por descenso de gradiente (módulos~4 y~7) para que el modelo prediga bien el token siguiente. La justificación de que este procedimiento produzca vectores con significado es la \emph{hipótesis distribucional}: las palabras que aparecen en contextos parecidos tienden a tener significados parecidos~\cite{harris}. Firth la resumió en la frase «you shall know a word by the company it keeps»~\cite{firth}. Dos tokens que aparecen en contextos similares reciben gradientes similares y acaban con embeddings próximos en el sentido del coseno~\cite[cap.~5]{jm3}.

Antes de los Transformers, los embeddings se entrenaban por separado con modelos como word2vec~\cite{mikolov-a}, que aprende un vector por palabra prediciendo las palabras vecinas en una ventana. En un LLM, en cambio, la matriz $E$ se aprende junto con el resto de los parámetros.

\section{La capa de salida y el atado de pesos}

Al final del modelo, el vector $h\in\R^{d}$ de la última posición se compara con un vector por cada token del vocabulario. Si $U\in\R^{K\times d}$ es la matriz de salida, los logits son
\[
z=Uh\in\R^{K},\qquad z_i=U_{i,:}\cdot h,
\]
y la softmax del módulo~1 (tema~4) los convierte en una distribución sobre el token siguiente. Cada logit es, por tanto, un producto escalar entre $h$ y un vector asociado al token.

\begin{definicion}[Atado de pesos]
Se dice que un modelo \emph{ata} los pesos de entrada y salida (\emph{weight tying}) si usa $U=E$: la misma matriz sirve para leer los embeddings de entrada y para calcular los logits.
\end{definicion}

Press y Wolf~\cite{press} mostraron que esta elección mejora los modelos de lenguaje recurrentes, y el Transformer original la adopta~\cite[§3.4]{vaswani}. Ahorra $Kd$ parámetros, que en GPT-2 son $38\,597\,376$, y da a ambos usos una interpretación común: el logit del token $i$ es alto cuando $h$ apunta en la dirección del embedding de $i$.

\section{Analogías y aritmética de vectores}

Sea $v_w$ el embedding de la palabra $w$. Si un rasgo como el género corresponde aproximadamente a una dirección del espacio, la diferencia $v_{\text{mujer}}-v_{\text{hombre}}$ debería parecerse a $v_{\text{reina}}-v_{\text{rey}}$. Mikolov, Yih y Zweig~\cite{mikolov-b} observaron este fenómeno en embeddings aprendidos y propusieron resolver analogías con el método siguiente.

\begin{definicion}[Analogía por desplazamiento de vectores]\label{def:analogia}
Dada la pregunta «$a$ es a $b$ como $?$ es a $c$», se calcula $q=v_a-v_b+v_c$ y se responde con
\[
\hat w=\operatorname*{arg\,max}_{w\in\V\setminus\{a,b,c\}}\cos(v_w,q).
\]
\end{definicion}

\begin{proposicion}\label{prop:paralelogramo}
Si $v_a-v_b=v_{w^{*}}-v_c$, entonces $q=v_{w^{*}}$ y $\cos(v_{w^{*}},q)=1$, de modo que $w^{*}$ es una respuesta de máxima similitud.
\end{proposicion}

\begin{proof}
$q=v_a-v_b+v_c=(v_{w^{*}}-v_c)+v_c=v_{w^{*}}$, y $\cos(u,u)=1$ para $u\neq 0$; por el \cref{teo:cs}, ningún candidato supera este valor.
\end{proof}

\begin{ejemplo}
Con los vectores del \cref{ej:juguete}, $v_{\text{rey}}-v_{\text{hombre}}+v_{\text{mujer}}=(2,1,2,0)-(0,1,2,0)+(0,-1,2,0)=(2,-1,2,0)=v_{\text{reina}}$. La resta elimina «masculino» y «persona»; la suma añade «femenino» y «persona»; la realeza se conserva.
\end{ejemplo}

\begin{observacion}[Limitaciones]
Con embeddings reales, la igualdad de la \cref{prop:paralelogramo} solo se cumple de forma aproximada, y $q$ suele quedar más cerca de $v_a$ que de la respuesta esperada; la exclusión de $a$, $b$ y $c$ en la \cref{def:analogia} es esencial para que el método «funcione». Nissim, van Noord y van der Goot~\cite{nissim} analizan cómo esta exclusión distorsiona las conclusiones que se extraen de las analogías. En el cuaderno de prácticas del módulo, con 40 palabras en 16 dimensiones y ruido, la consulta rey $-$ hombre $+$ mujer da «reina» con coseno $0{,}924$, pero sin la exclusión «rey» queda segunda con $0{,}835$. Además, los embeddings heredan los sesgos de los textos de entrenamiento: Bolukbasi y coautores~\cite{bolukbasi} encontraron analogías que reproducen estereotipos de género. La idea de que los conceptos son direcciones reaparece en la interpretabilidad de los LLM (módulo~23).
\end{observacion}

\section{Adelanto: el producto escalar en la atención}

La atención (módulo~13) es un mecanismo de búsqueda por similitud. Cada token emite una consulta $q\in\R^{d_k}$, la compara por producto escalar con las claves $k_1,\dots,k_n$ de los tokens anteriores y reparte su atención con una softmax:
\[
s_j=\frac{q\cdot k_j}{\sqrt{d_k}},\qquad \alpha=\softmax(s_1,\dots,s_n).
\]
Para todos los tokens a la vez, las puntuaciones forman la matriz $QK^{\top}/\sqrt{d_k}$ (\cref{def:gram}).

\begin{ejemplo}\label{ej:atencion}
Sean $d_k=2$, $q=(1,0)$ y las claves $k_1=(1,1)$, $k_2=(0,2)$, $k_3=(-1,0)$. Los productos escalares son $1$, $0$ y $-1$; divididos por $\sqrt2$ dan $(0{,}7071;\,0;\,-0{,}7071)$, y la softmax da $\alpha\approx(0{,}5760;\,0{,}2840;\,0{,}1400)$. El token cuya clave está más alineada con la consulta recibe más atención.
\end{ejemplo}

El factor $1/\sqrt{d_k}$ tiene una justificación estadística~\cite[§3.2.1]{vaswani}: si las componentes de $q$ y $k$ son variables independientes de media~0 y varianza~1, entonces $\Var(q\cdot k)=d_k$ (ejercicio~\ref{eje:5-atencion}). Sin el factor, para $d_k$ grande las puntuaciones tendrían valores muy dispares, la softmax quedaría casi concentrada en un solo token y los gradientes serían casi nulos. Dividir por $\sqrt{d_k}$ devuelve la varianza a~1.

\begin{lecturas}
\item Jurafsky y Martin~\cite{jm3}, cap.~5 «Embeddings».
\item CS224N~\cite{cs224n}, clase~2 «Word Vectors» y tarea~1.
\item Zhang \emph{et al.}~\cite{d2l}, §15.7 «Word Similarity and Analogy».
\item Mikolov, Yih y Zweig~\cite{mikolov-b}; Vaswani \emph{et al.}~\cite{vaswani}, §3.2.1 y §3.4.
\end{lecturas}

\section*{Ejercicios del tema 5}
\addcontentsline{toc}{section}{Ejercicios del tema 5}

\begin{ejercicio}{1}{eje:5-tamano}
GPT-2 tiene $K=50\,257$ y, en su versión más pequeña, $d=768$. ¿Qué tamaño tiene $E$? ¿Cuántos parámetros tiene y cuánto ocupa con 4 bytes por número? Repita el cálculo con $d=1600$, la dimensión de la versión más grande.
\end{ejercicio}

\begin{ejercicio}{1}{eje:5-onehot}
Explique, a partir de la definición de producto de matrices, por qué $e_i^{\top}E$ es la fila $i$ de $E$, y estime cuántas veces más operaciones requiere el producto que la lectura directa de la fila en GPT-2.
\end{ejercicio}

\begin{ejercicio}{1}{eje:5-busqueda}
Se desea encontrar las palabras más similares a una consulta $q$ en un vocabulario muy grande. (a) Explique por qué, si se normalizan una vez todas las filas de $E$ para formar $\hat E$, los $K$ cosenos se obtienen con un único producto matriz-vector. (b) Cuente las multiplicaciones necesarias con $K=50\,257$ y $d=768$. (c) ¿Cómo se procesa a la vez un lote de $B$ consultas?
\fuente{adaptado de Zhang \emph{et al.}~\cite{d2l}, §15.7, ejercicio~2.}
\end{ejercicio}

\begin{ejercicio}{2}{eje:5-analogias}
Con los vectores del \cref{ej:juguete}: (a) calcule $\cos(\text{hombre},\text{mujer})$; (b) compruebe que $v_{\text{rey}}-v_{\text{hombre}}+v_{\text{mujer}}=v_{\text{reina}}$. (c) Con $v_{\text{España}}=(1,0,1,0)$, $v_{\text{Madrid}}=(0,1,1,0)$, $v_{\text{Francia}}=(1,0,0,1)$, $v_{\text{París}}=(0,1,0,1)$ y $v_{\text{Lyon}}=(0;\,0{,}5;\,0;\,1)$, calcule $q=v_{\text{Madrid}}-v_{\text{España}}+v_{\text{Francia}}$ y su coseno con París y con Lyon.
\end{ejercicio}

\begin{ejercicio}{2}{eje:5-secuencia}
Con la matriz $E$ del \cref{ej:embedding}: (a) calcule $e_2^{\top}E$ columna a columna; (b) escriba la matriz $O$ de la \cref{prop:onehot} para la secuencia $(2,0,2)$ y calcule $X=OE$ indicando su tamaño.
\end{ejercicio}

\begin{ejercicio}{2}{eje:5-logits}
Un modelo con pesos atados usa la matriz $E$ del \cref{ej:embedding} como matriz de salida. El vector final es $h=(1,0,-1)$. Calcule los logits $z=Eh$ y la distribución $\softmax(z)$. ¿Qué token es el más probable y por qué, en términos de ángulos?
\end{ejercicio}

\begin{ejercicio}{3}{eje:5-atencion}
(a) Reproduzca los cálculos del \cref{ej:atencion}. (b) Sean $q,k\in\R^{d}$ vectores aleatorios cuyas $2d$ componentes son independientes, de media~0 y varianza~1. Demuestre que $\E[q\cdot k]=0$ y $\Var(q\cdot k)=d$. (c) Con $d=64$, compare la softmax de las puntuaciones típicas $(8,0,-8)$, del orden de una desviación típica, con la de $(1,0,-1)$, que se obtiene al dividir por $\sqrt{64}$.
\end{ejercicio}


\appendix
% =====================================================================
\chapter{Soluciones de los ejercicios}
\label{ap:soluciones}
% =====================================================================

Cada solución se presenta en cinco partes: datos, planteamiento, desarrollo, resultado y comprobación. En los pasos intermedios se redondea a cuatro cifras significativas salvo que se indique otra cosa; los resultados exactos se dan como fracción o radical cuando existe.

% ---------------------------------------------------------------------
\section{Tema 1}

\begin{solucion}{eje:1-combinacion}
\datos $u=(1,2)$, $v=(3,-1)$, objetivo $(5,3)$.
\planteamiento igualar componente a componente $au+bv=(a+3b,\,2a-b)$ con $(5,3)$ y resolver el sistema.
\pasos
\begin{enumerate}[label=\arabic*.]
  \item Sistema: $a+3b=5$ y $2a-b=3$.
  \item De la segunda ecuación, $b=2a-3$. Sustituyendo en la primera: $a+3(2a-3)=5$, es decir, $7a-9=5$, luego $a=2$.
  \item $b=2\cdot 2-3=1$.
\end{enumerate}
\resultado \textbf{$a=2$, $b=1$: $(5,3)=2u+v$.} Geométricamente, $(5,3)$ es la diagonal del paralelogramo de lados $2u$ y $v$.
\comprobacion $2(1,2)+(3,-1)=(2+3,\,4-1)=(5,3)$, que coincide con el \cref{ej:operaciones}.
\end{solucion}

\begin{solucion}{eje:1-recuentos}
\datos $w\in\R^{n}$ vector de recuentos; diccionario (el, gato, come, pez); documento «el gato come el pez».
\planteamiento interpretar cada producto escalar y cada componente según la definición de $w$.
\pasos
\begin{enumerate}[label=(\alph*)]
  \item $\unos\cdot w=\sum_iw_i$ suma los recuentos de todas las palabras del diccionario: es el \textbf{número total de palabras} del documento.
  \item $w_{282}=0$ significa que la palabra 282 del diccionario \textbf{no aparece} en el documento.
  \item $h_i=w_i/(\unos\cdot w)$, es decir, $h=\dfrac{1}{\unos\cdot w}\,w$, definido si el documento tiene al menos una palabra.
  \item «el» aparece dos veces y las demás una: $w=(2,1,1,1)$; $\unos\cdot w=2+1+1+1=5$; $h=\tfrac15(2,1,1,1)=(0{,}4;\,0{,}2;\,0{,}2;\,0{,}2)$.
\end{enumerate}
\resultado \textbf{(a) número total de palabras; (b) la palabra 282 no aparece; (c) $h=w/(\unos\cdot w)$; (d) $w=(2,1,1,1)$, $\unos\cdot w=5$, $h=(0{,}4;\,0{,}2;\,0{,}2;\,0{,}2)$.}
\comprobacion $\unos\cdot h=0{,}4+0{,}2+0{,}2+0{,}2=1$: el histograma es una distribución de probabilidad sobre el diccionario.
\end{solucion}

\begin{solucion}{eje:1-normas}
\datos $x=(3,-4,12)$.
\planteamiento aplicar las definiciones de $\norm{\cdot}_1$, $\norm{\cdot}$ y $\norm{\cdot}_\infty$; después dividir por la norma euclídea.
\pasos
\begin{enumerate}[label=\arabic*.]
  \item $\norm{x}_1=\abs{3}+\abs{-4}+\abs{12}=3+4+12=19$.
  \item $\norm{x}=\sqrt{9+16+144}=\sqrt{169}=13$.
  \item $\norm{x}_\infty=\max\{3,4,12\}=12$.
  \item Orden: $12\le 13\le 19$.
  \item $\hat x=x/13=(3/13,\,-4/13,\,12/13)\approx(0{,}2308;\,-0{,}3077;\,0{,}9231)$.
\end{enumerate}
\resultado \textbf{$\norm{x}_1=19$, $\norm{x}=13$, $\norm{x}_\infty=12$; $\hat x=\frac{1}{13}(3,-4,12)$.}
\comprobacion $\norm{\hat x}^{2}=(9+16+144)/169=1$.
\end{solucion}

\begin{solucion}{eje:1-cs}
\datos $u,v\in\R^{d}$ con $v\neq 0$ (si $v=0$ ambos miembros valen~0).
\planteamiento un polinomio de segundo grado con coeficiente principal positivo que no toma valores negativos tiene discriminante $\le 0$.
\pasos
\begin{enumerate}[label=\arabic*.]
  \item $f(t)=\norm{u-tv}^{2}\ge 0$ para todo $t\in\R$, por ser una norma al cuadrado.
  \item Por la linealidad y la simetría del producto escalar, $f(t)=(u-tv)\cdot(u-tv)=\norm{v}^{2}t^{2}-2(u\cdot v)\,t+\norm{u}^{2}$.
  \item Es un polinomio $\alpha t^{2}+\beta t+\gamma$ con $\alpha=\norm{v}^{2}>0$, $\beta=-2(u\cdot v)$ y $\gamma=\norm{u}^{2}$. Si su discriminante fuera positivo, tendría dos raíces reales distintas y sería negativo entre ellas; como $f\ge 0$, debe ser $\beta^{2}-4\alpha\gamma\le 0$.
  \item Sustituyendo: $4(u\cdot v)^{2}-4\norm{u}^{2}\norm{v}^{2}\le 0$, es decir, $(u\cdot v)^{2}\le\norm{u}^{2}\norm{v}^{2}$.
  \item Tomando raíces cuadradas (ambos miembros son $\ge 0$): $\abs{u\cdot v}\le\norm{u}\norm{v}$.
  \item Si $u,v\neq 0$, dividiendo por $\norm{u}\norm{v}>0$: $\abs{u\cdot v}/(\norm{u}\norm{v})\le 1$.
\end{enumerate}
\resultado \textbf{$\abs{u\cdot v}\le\norm{u}\norm{v}$ y, por tanto, $-1\le u\cdot v/(\norm{u}\norm{v})\le 1$.}
\comprobacion con $u=(1,2)$ y $v=(3,4)$: $\abs{u\cdot v}=11$ y $\norm{u}\norm{v}=\sqrt5\cdot 5\approx 11{,}18$. La desigualdad se cumple con poco margen porque los vectores son casi paralelos.
\end{solucion}

\begin{solucion}{eje:1-paralelogramo}
\datos $u,v\in\R^{d}$; en (c), $u=(1,0)$ y $v=(0,1)$.
\planteamiento desarrollar $\norm{u\pm v}^{2}$ con la linealidad del producto escalar, y evaluar ambos miembros con la norma $\ell_1$.
\pasos
\begin{enumerate}[label=\arabic*.]
  \item $\norm{u+v}^{2}=\norm{u}^{2}+2\,u\cdot v+\norm{v}^{2}$ y $\norm{u-v}^{2}=\norm{u}^{2}-2\,u\cdot v+\norm{v}^{2}$.
  \item (a) Sumando: $\norm{u+v}^{2}+\norm{u-v}^{2}=2\norm{u}^{2}+2\norm{v}^{2}$.
  \item (b) Restando: $\norm{u+v}^{2}-\norm{u-v}^{2}=4\,u\cdot v$, de donde se obtiene la identidad de polarización.
  \item (c) $u+v=(1,1)$ y $u-v=(1,-1)$, luego $\norm{u+v}_1=\norm{u-v}_1=2$ y el primer miembro vale $2^{2}+2^{2}=8$.
  \item El segundo miembro vale $2\norm{u}_1^{2}+2\norm{v}_1^{2}=2+2=4\neq 8$.
  \item Los pasos 1 y 2 solo usan la simetría y la linealidad, válidas para todo producto interior. Si la norma $\ell_1$ fuera $\sqrt{\ip{u}{u}}$ para algún producto interior, cumpliría la ley del paralelogramo, lo que contradice el paso~5.
\end{enumerate}
\resultado \textbf{Se cumplen la ley del paralelogramo y la identidad de polarización; la norma $\ell_1$ no cumple la primera ($8\neq 4$) y, por tanto, no procede de ningún producto interior.}
\comprobacion con la norma euclídea y los mismos vectores: $\norm{u+v}^{2}+\norm{u-v}^{2}=2+2=4=2\cdot 1+2\cdot 1$; y la polarización da $\tfrac14(2-2)=0=u\cdot v$.
\end{solucion}

% ---------------------------------------------------------------------
\section{Tema 2}

\begin{solucion}{eje:2-escala}
\datos $a=(1,2)$, $b=(2,4)$, $c=(10,20)$.
\planteamiento $u\cdot v=\sum_iu_iv_i$ y $\cos(u,v)=u\cdot v/(\norm{u}\norm{v})$.
\pasos
\begin{enumerate}[label=\arabic*.]
  \item $a\cdot b=1\cdot 2+2\cdot 4=10$ y $a\cdot c=1\cdot 10+2\cdot 20=50$.
  \item $\norm{a}=\sqrt5$, $\norm{b}=\sqrt{20}$, $\norm{c}=\sqrt{500}$.
  \item $\cos(a,b)=10/\sqrt{5\cdot 20}=10/10=1$ y $\cos(a,c)=50/\sqrt{5\cdot 500}=50/50=1$.
\end{enumerate}
\resultado \textbf{$a\cdot b=10$, $a\cdot c=50$; ambos cosenos valen~1.} El producto escalar crece con la longitud de los vectores; la similitud coseno solo mide la dirección, que es la misma en los tres.
\comprobacion $c=10a$, luego $a\cdot c=10\,(a\cdot a)=10\cdot 5=50$; y, por la \cref{prop:coseno}(ii), vectores múltiplos positivos tienen coseno~1.
\end{solucion}

\begin{solucion}{eje:2-vf}
\planteamiento buscar un contraejemplo o aplicar la definición.
\pasos
\begin{enumerate}[label=(\alph*)]
  \item \textbf{Falsa.} $(1,0)\cdot(0,1)=0$ y ninguno es nulo: $u\cdot v=0$ significa que son ortogonales.
  \item \textbf{Verdadera.} $\cos(u,u)=u\cdot u/(\norm{u}\norm{u})=\norm{u}^{2}/\norm{u}^{2}=1$ si $u\neq 0$.
  \item \textbf{Falsa.} $\cos((1,2),(2,4))=1$ y los vectores son distintos. Por la \cref{prop:coseno}(ii), $\cos(u,v)=1$ solo implica $u=cv$ con $c>0$.
\end{enumerate}
\resultado \textbf{(a) falsa; (b) verdadera; (c) falsa.}
\comprobacion en (b), con $u=(3,4)$: $25/(5\cdot 5)=1$.
\end{solucion}

\begin{solucion}{eje:2-angulo}
\datos $u=(2,-1,2)$, $v=(1,2,2)$.
\planteamiento producto escalar, normas, coseno y $\theta=\arccos(\cos(u,v))$ (\cref{def:angulo}).
\pasos
\begin{enumerate}[label=\arabic*.]
  \item $u\cdot v=2\cdot 1+(-1)\cdot 2+2\cdot 2=2-2+4=4$.
  \item $\norm{u}=\sqrt{4+1+4}=3$ y $\norm{v}=\sqrt{1+4+4}=3$.
  \item $\cos(u,v)=4/(3\cdot 3)=4/9\approx 0{,}4444$.
  \item $\theta=\arccos(0{,}4444)\approx 63{,}61^\circ$.
\end{enumerate}
\resultado \textbf{$u\cdot v=4$; $\norm{u}=\norm{v}=3$; $\cos(u,v)=4/9\approx 0{,}4444$; $\theta\approx 63{,}6^\circ$.}
\comprobacion $\cos 60^\circ=0{,}5$ y $\cos 90^\circ=0$; un coseno de $0{,}44$ corresponde a un ángulo entre ambos, más próximo a $60^\circ$.
\end{solucion}

\begin{solucion}{eje:2-vecinos}
\datos $a=(1,1)$, $b=(3,3)$, $c=(1,-1)$.
\planteamiento comparar la distancia $\norm{a-x}$ con la similitud $\cos(a,x)$ para $x=b,c$.
\pasos
\begin{enumerate}[label=\arabic*.]
  \item $a-b=(-2,-2)$ y $\norm{a-b}=\sqrt8\approx 2{,}828$.
  \item $a-c=(0,2)$ y $\norm{a-c}=2$.
  \item $a\cdot b=6$, $\norm{a}=\sqrt2$, $\norm{b}=\sqrt{18}$; $\cos(a,b)=6/\sqrt{36}=1$.
  \item $a\cdot c=1-1=0$, luego $\cos(a,c)=0$.
  \item Normalizando: $\hat a=\hat b=(1/\sqrt2,1/\sqrt2)$ y $\hat c=(1/\sqrt2,-1/\sqrt2)$. Entonces $\norm{\hat a-\hat b}=0$ y $\norm{\hat a-\hat c}=\norm{(0,\sqrt2)}=\sqrt2$.
\end{enumerate}
\resultado \textbf{$\norm{a-b}\approx 2{,}828$, $\norm{a-c}=2$, $\cos(a,b)=1$, $\cos(a,c)=0$.} La distancia elige $c$ y el coseno elige $b$: $b$ tiene la dirección de $a$, pero es tres veces más largo. Tras normalizar, ambas medidas eligen $b$.
\comprobacion la \cref{prop:coseno-distancia} da $\norm{\hat a-\hat c}^{2}=2-2\cdot 0=2$, que coincide con $(\sqrt2)^{2}$.
\end{solucion}

\begin{solucion}{eje:2-normalizar}
\datos $x=(3,4)$; $r=(1,0)$.
\planteamiento $\hat x=x/\norm{x}$; después evaluar ambos miembros de la identidad.
\pasos
\begin{enumerate}[label=\arabic*.]
  \item $\norm{x}=\sqrt{9+16}=5$ y $p=\hat x=(0{,}6;\,0{,}8)$.
  \item $p-r=(-0{,}4;\,0{,}8)$ y $\norm{p-r}^{2}=0{,}16+0{,}64=0{,}80$.
  \item Como $\norm{p}=\norm{r}=1$, $\cos(p,r)=p\cdot r=0{,}6$.
  \item $2-2\cos(p,r)=2-1{,}2=0{,}80$.
\end{enumerate}
\resultado \textbf{$\hat x=(0{,}6;\,0{,}8)$, y ambos miembros valen $0{,}80$.}
\comprobacion $\norm{p}^{2}=0{,}36+0{,}64=1$: $p$ es unitario y la identidad es aplicable.
\end{solucion}

\begin{solucion}{eje:2-tres}
\datos vectores de $\R^{2}$ con todos los productos escalares por pares negativos.
\planteamiento producto escalar negativo equivale a ángulo mayor que $90^\circ$ (\cref{tab:signo}); tres ángulos de $120^\circ$ completan la vuelta.
\pasos
\begin{enumerate}[label=\arabic*.]
  \item Se toman $u=(1,0)$, $v=(-1/2,\sqrt3/2)$ y $w=(-1/2,-\sqrt3/2)$, que son unitarios: $\norm{v}^{2}=\tfrac14+\tfrac34=1$, y lo mismo $w$.
  \item $u\cdot v=-\tfrac12$; $u\cdot w=-\tfrac12$; $v\cdot w=\tfrac14-\tfrac34=-\tfrac12$.
  \item Como son unitarios, $\cos\theta=-\tfrac12$ en los tres pares, luego $\theta=120^\circ$.
  \item (b) Supóngase que cuatro vectores no nulos tienen los seis productos negativos. Se ordenan por su ángulo polar y se llaman $g_1,\dots,g_4$ los ángulos que separan cada vector del siguiente en sentido antihorario, de modo que $g_1+g_2+g_3+g_4=360^\circ$.
  \item El ángulo entre dos vectores consecutivos es $\min(g_i,360^\circ-g_i)\le g_i$, y debe ser mayor que $90^\circ$; luego $g_i>90^\circ$ para todo $i$ y la suma supera $360^\circ$, lo cual es imposible.
\end{enumerate}
\resultado \textbf{(a) Sí: tres vectores unitarios a $120^\circ$, con los tres productos iguales a $-1/2$. (b) No: cuatro vectores del plano no pueden tener los seis productos negativos.}
\comprobacion $u+v+w=(1-\tfrac12-\tfrac12,\ 0+\tfrac{\sqrt3}{2}-\tfrac{\sqrt3}{2})=(0,0)$: las tres flechas se equilibran, como tres fuerzas iguales a $120^\circ$. En $\R^{3}$, en cambio, los cuatro vértices de un tetraedro regular centrado en el origen sí tienen los seis productos negativos.
\end{solucion}

\begin{solucion}{eje:2-distancia}
\datos $\norm{u}=\norm{v}=1$.
\planteamiento desarrollar la norma al cuadrado mediante el producto escalar.
\pasos
\begin{enumerate}[label=\arabic*.]
  \item $\norm{u-v}^{2}=(u-v)\cdot(u-v)=\norm{u}^{2}-2\,u\cdot v+\norm{v}^{2}$, por la linealidad y la simetría.
  \item Con normas unitarias: $\norm{u-v}^{2}=2-2\,u\cdot v$.
  \item También $\cos(u,v)=u\cdot v/(1\cdot 1)=u\cdot v$, luego $\norm{u-v}^{2}=2-2\cos(u,v)$.
  \item La función $c\mapsto 2-2c$ es estrictamente decreciente. Si $\cos(q,a)>\cos(q,b)$, entonces $\norm{q-a}^{2}<\norm{q-b}^{2}$ y, tomando raíces, $\norm{q-a}<\norm{q-b}$.
  \item Si $\cos(u,v)=-1$: $\norm{u-v}^{2}=4$ y $\norm{u-v}=2$.
\end{enumerate}
\resultado \textbf{$\norm{u-v}^{2}=2-2\cos(u,v)$; entre vectores unitarios, el más cercano por distancia es el de mayor coseno. Con $\cos=-1$, la distancia es~2 (puntos antipodales de la esfera unidad).}
\comprobacion el ejercicio~\ref{eje:2-normalizar} da $0{,}80=2-2\cdot 0{,}6$; con $\cos=1$ la distancia es~0, como debe ser para $u=v$.
\end{solucion}

% ---------------------------------------------------------------------
\section{Tema 3}

\begin{solucion}{eje:3-lineal}
\datos $f(x)=(x_1+x_2,3x_1)$, $g(x)=(x_1+1,x_2)$, $h(x)=(x_1x_2,x_2)$; $W\in\R^{5\times 3}$.
\planteamiento una aplicación lineal cumple $T(0)=0$ y $T(cx)=cT(x)$; si es lineal, su matriz tiene por columnas $T(e_1)$ y $T(e_2)$ (\cref{teo:representacion}).
\pasos
\begin{enumerate}[label=\arabic*.]
  \item $f$ es de la forma $Ax$ con $A=\begin{pmatrix}1&1\\3&0\end{pmatrix}$: $Ax=(x_1+x_2,\,3x_1)$. Por la \cref{prop:matriz-lineal}, es lineal. Columnas: $f(e_1)=(1,3)$ y $f(e_2)=(1,0)$.
  \item $g(0)=(1,0)\neq 0$, luego $g$ no es lineal.
  \item $h(2,2)=(4,2)$ y $2h(1,1)=2(1,1)=(2,2)$; como difieren, $h$ no es lineal.
  \item (b) $W$ actúa sobre vectores de $\R^{3}$ (tres columnas) y produce vectores de $\R^{5}$ (cinco filas).
\end{enumerate}
\resultado \textbf{Solo $f$ es lineal, con matriz $\begin{pmatrix}1&1\\3&0\end{pmatrix}$. $W$ transforma $\R^{3}$ en $\R^{5}$.}
\comprobacion $f(1,1)=(2,3)$ y $A(1,1)=(1+1,\,3+0)=(2,3)$.
\end{solucion}

\begin{solucion}{eje:3-dos-lecturas}
\datos $M=\begin{pmatrix}1&0&2\\2&1&0\\0&3&1\end{pmatrix}$, $x=(2,-1,3)$.
\planteamiento las dos lecturas de la \cref{prop:filas-columnas}.
\pasos
\begin{enumerate}[label=\arabic*.]
  \item Fila 1: $1\cdot 2+0\cdot(-1)+2\cdot 3=8$.
  \item Fila 2: $2\cdot 2+1\cdot(-1)+0\cdot 3=3$.
  \item Fila 3: $0\cdot 2+3\cdot(-1)+1\cdot 3=0$.
  \item Columnas: $2(1,2,0)+(-1)(0,1,3)+3(2,0,1)=(2,4,0)+(0,-1,-3)+(6,0,3)=(8,3,0)$.
\end{enumerate}
\resultado \textbf{$Mx=(8,3,0)$.} La tercera fila de $M$ es ortogonal a $x$, porque su producto escalar es~0.
\comprobacion las dos lecturas dan el mismo vector, como garantiza la \cref{prop:filas-columnas}.
\end{solucion}

\begin{solucion}{eje:3-base}
\datos $T(e_1)=(1,2)$, $T(e_2)=(-1,1)$; $x=(3,-2)$.
\planteamiento por el \cref{teo:representacion}, las columnas de la matriz son $T(e_1)$ y $T(e_2)$.
\pasos
\begin{enumerate}[label=\arabic*.]
  \item $A=\begin{pmatrix}1&-1\\2&1\end{pmatrix}$.
  \item Con la matriz: $Ax=(1\cdot 3+(-1)(-2),\ 2\cdot 3+1\cdot(-2))=(5,4)$.
  \item Con la linealidad: $T(3e_1-2e_2)=3T(e_1)-2T(e_2)=(3,6)+(2,-2)=(5,4)$.
\end{enumerate}
\resultado \textbf{$A=\begin{pmatrix}1&-1\\2&1\end{pmatrix}$ y $T(3,-2)=(5,4)$.}
\comprobacion $Ae_1=(1,2)$ y $Ae_2=(-1,1)$, que son los datos.
\end{solucion}

\begin{solucion}{eje:3-rotacion}
\datos $\theta=60^\circ$, $\cos 60^\circ=\tfrac12$, $\sin 60^\circ=\tfrac{\sqrt3}{2}$.
\planteamiento sustituir en la matriz $R_\theta$ del \cref{tab:transformaciones} y multiplicar.
\pasos
\begin{enumerate}[label=\arabic*.]
  \item $R_{60^\circ}=\begin{pmatrix}1/2&-\sqrt3/2\\ \sqrt3/2&1/2\end{pmatrix}$.
  \item $R_{60^\circ}(2,0)=(1,\sqrt3)$, cuya norma es $\sqrt{1+3}=2=\norm{(2,0)}$.
  \item $R_{60^\circ}(1,1)=\bigl(\tfrac{1-\sqrt3}{2},\tfrac{\sqrt3+1}{2}\bigr)\approx(-0{,}3660;\,1{,}366)$.
  \item Norma al cuadrado: $\tfrac14\bigl[(1-\sqrt3)^{2}+(1+\sqrt3)^{2}\bigr]=\tfrac14\bigl[(4-2\sqrt3)+(4+2\sqrt3)\bigr]=2$, luego la norma es $\sqrt2=\norm{(1,1)}$.
\end{enumerate}
\resultado \textbf{$R_{60^\circ}(2,0)=(1,\sqrt3)$ y $R_{60^\circ}(1,1)\approx(-0{,}3660;\,1{,}366)$; las normas ($2$ y $\sqrt2$) se conservan.}
\comprobacion el vector $(1,\sqrt3)$ forma con el eje $x$ un ángulo de $\arctan\sqrt3=60^\circ$, como se esperaba.
\end{solucion}

\begin{solucion}{eje:3-markov}
\datos $A=\begin{pmatrix}0{,}8&0{,}3\\0{,}2&0{,}7\end{pmatrix}$, $u_0=(1,0)$, $s=(0{,}6;\,0{,}4)$.
\planteamiento aplicar $A$ repetidamente; para (c), intercambiar el orden de las sumas.
\pasos
\begin{enumerate}[label=\arabic*.]
  \item $u_1=Au_0=(0{,}8;\,0{,}2)$, la primera columna de $A$.
  \item $u_2=Au_1=(0{,}8\cdot 0{,}8+0{,}3\cdot 0{,}2;\ 0{,}2\cdot 0{,}8+0{,}7\cdot 0{,}2)=(0{,}64+0{,}06;\ 0{,}16+0{,}14)=(0{,}70;\,0{,}30)$.
  \item $u_3=Au_2=(0{,}8\cdot 0{,}7+0{,}3\cdot 0{,}3;\ 0{,}2\cdot 0{,}7+0{,}7\cdot 0{,}3)=(0{,}56+0{,}09;\ 0{,}14+0{,}21)=(0{,}65;\,0{,}35)$.
  \item Las componentes de $u_0,\dots,u_3$ son no negativas y suman~1.
  \item (b) $As=(0{,}48+0{,}12;\ 0{,}12+0{,}28)=(0{,}6;\,0{,}4)=s$.
  \item (c) $\unos^{\top}(Au)=\sum_{i}\sum_{j}A_{ij}u_j=\sum_ju_j\Bigl(\sum_iA_{ij}\Bigr)=\sum_ju_j\cdot 1=\unos^{\top}u$.
\end{enumerate}
\resultado \textbf{$u_1=(0{,}8;\,0{,}2)$, $u_2=(0{,}7;\,0{,}3)$, $u_3=(0{,}65;\,0{,}35)$, todos con suma~1; $As=s$; una matriz con columnas de suma~1 conserva la suma de las componentes.}
\comprobacion los $u_k$ se acercan a $s$: la primera componente pasa de $1$ a $0{,}8$, $0{,}7$ y $0{,}65$, y la distancia a $0{,}6$ se reduce a la mitad en cada paso ($0{,}4$; $0{,}2$; $0{,}1$; $0{,}05$).
\end{solucion}

\begin{solucion}{eje:3-rectas}
\datos $T\colon\R^{n}\to\R^{m}$ lineal; $p,d\in\R^{n}$; $b\neq 0$.
\planteamiento aplicar la \cref{def:lineal} a combinaciones concretas.
\pasos
\begin{enumerate}[label=\arabic*.]
  \item (a) Por linealidad, $T(p+td)=T(p)+t\,T(d)$ para todo $t$. Si $T(d)\neq 0$, la imagen es la recta que pasa por $T(p)$ con dirección $T(d)$; si $T(d)=0$, la imagen es el punto $T(p)$.
  \item (b) Con $a=b=\tfrac12$ en la \cref{def:lineal}: $T\bigl(\tfrac12u+\tfrac12v\bigr)=\tfrac12T(u)+\tfrac12T(v)$.
  \item (c) Sea $S(x)=x+b$. La imagen de la recta es $\{(p+b)+td\}$, otra recta, y $S\bigl(\tfrac12(u+v)\bigr)=\tfrac12(u+v)+b=\tfrac12\bigl((u+b)+(v+b)\bigr)$.
  \item Sin embargo, $S(0)=b\neq 0$, mientras que toda aplicación lineal fija el origen. Luego $S$ no es lineal.
\end{enumerate}
\resultado \textbf{Una aplicación lineal transforma rectas en rectas (o en un punto, si $T(d)=0$) y conserva puntos medios; las traslaciones tienen ambas propiedades sin ser lineales.} Las aplicaciones de la forma $x\mapsto Ax+b$ se llaman \emph{afines}~\cite[§8.1]{vmls}; una capa de red neuronal con sesgo es afín (módulo~9).
\comprobacion la proyección sobre el eje $x$ del \cref{tab:transformaciones} envía la recta vertical $\{(1,0)+t(0,1)\}$ al punto $(1,0)$, porque $T(0,1)=(0,0)$.
\end{solucion}

% ---------------------------------------------------------------------
\section{Tema 4}

\begin{solucion}{eje:4-formas}
\datos tamaños de las matrices y $u,v\in\R^{3}$.
\planteamiento $(m\times n)(n\times p)=m\times p$, solo si coinciden las dimensiones interiores.
\pasos
\begin{enumerate}[label=(\alph*)]
  \item $(3\times 4)(4\times 2)=3\times 2$. El producto $(2\times 3)(2\times 3)$ no está definido, porque $3\neq 2$.
  \item $u^{\top}v$ es $(1\times 3)(3\times 1)=1\times 1$, un número (el producto escalar); $uv^{\top}$ es $(3\times 1)(1\times 3)=3\times 3$, un producto exterior.
  \item \textbf{Falso.} En el \cref{ej:intercambio}, $AB\neq BA$.
\end{enumerate}
\resultado \textbf{(a) $3\times 2$; no. (b) $1\times 1$ y $3\times 3$. (c) Falso.}
\comprobacion con $u=v=(1,1,1)$: $u^{\top}v=3$, y $uv^{\top}$ es la matriz $3\times 3$ de unos, cuya traza es también~3.
\end{solucion}

\begin{solucion}{eje:4-rs}
\datos $R=\begin{pmatrix}0&-1\\1&0\end{pmatrix}$, $S=\begin{pmatrix}2&0\\0&0{,}5\end{pmatrix}$, $x=(3,1)$.
\planteamiento entrada $(i,j)$ del producto $=$ fila $i$ del primer factor por columna $j$ del segundo.
\pasos
\begin{enumerate}[label=\arabic*.]
  \item $RS$: fila 1, $(0\cdot 2+(-1)\cdot 0;\ 0\cdot 0+(-1)\cdot 0{,}5)=(0;\,-0{,}5)$; fila 2, $(1\cdot 2+0;\ 0+0)=(2;\,0)$.
  \item $SR$: fila 1, $(2\cdot 0+0\cdot 1;\ 2\cdot(-1)+0\cdot 0)=(0;\,-2)$; fila 2, $(0\cdot 0+0{,}5\cdot 1;\ 0+0)=(0{,}5;\,0)$.
  \[
  RS=\begin{pmatrix}0&-0{,}5\\2&0\end{pmatrix},\qquad SR=\begin{pmatrix}0&-2\\0{,}5&0\end{pmatrix}.
  \]
  \item $RSx=(0\cdot 3-0{,}5\cdot 1;\ 2\cdot 3+0)=(-0{,}5;\,6)$. Por pasos: $Sx=(6;\,0{,}5)$ y $R(6;\,0{,}5)=(-0{,}5;\,6)$.
  \item $SRx=(0\cdot 3-2\cdot 1;\ 0{,}5\cdot 3+0)=(-2;\,1{,}5)$. Por pasos: $Rx=(-1,3)$ y $S(-1,3)=(-2;\,1{,}5)$.
\end{enumerate}
\resultado \textbf{$RS=\begin{pmatrix}0&-0{,}5\\2&0\end{pmatrix}$, $SR=\begin{pmatrix}0&-2\\0{,}5&0\end{pmatrix}$; $RSx=(-0{,}5;\,6)$ y $SRx=(-2;\,1{,}5)$.} Los resultados coinciden con la aplicación sucesiva, como afirma el \cref{teo:composicion}.
\comprobacion $R$ conserva normas: $\norm{Sx}^{2}=36+0{,}25=36{,}25$ y $\norm{R(Sx)}^{2}=0{,}25+36=36{,}25$.
\end{solucion}

\begin{solucion}{eje:4-ab-ba}
\datos $A\in\R^{2\times 3}$ y $B\in\R^{3\times 2}$ del enunciado.
\planteamiento $(2\times 3)(3\times 2)=2\times 2$ y $(3\times 2)(2\times 3)=3\times 3$.
\pasos
\begin{enumerate}[label=\arabic*.]
  \item $(AB)_{11}=1\cdot 1+0\cdot 0+2\cdot 3=7$; $(AB)_{12}=1\cdot 2+0\cdot 1+2\cdot 0=2$.
  \item $(AB)_{21}=0\cdot 1+1\cdot 0+(-1)\cdot 3=-3$; $(AB)_{22}=0\cdot 2+1\cdot 1+(-1)\cdot 0=1$.
  \item Filas de $BA$ (\cref{prop:cuatro}(ii)): $(1,2)A=(1,2,0)$; $(0,1)A=(0,1,-1)$; $(3,0)A=(3,0,6)$.
  \[
  AB=\begin{pmatrix}7&2\\-3&1\end{pmatrix},\qquad BA=\begin{pmatrix}1&2&0\\0&1&-1\\3&0&6\end{pmatrix}.
  \]
  \item $AA$ sería $(2\times 3)(2\times 3)$, que no está definido.
\end{enumerate}
\resultado \textbf{$AB$ es $2\times 2$ y $BA$ es $3\times 3$, con los valores indicados; $AA$ no existe.}
\comprobacion las trazas de $AB$ y de $BA$ coinciden siempre: $7+1=8$ y $1+1+6=8$.
\end{solucion}

\begin{solucion}{eje:4-distributiva}
\datos $A=\begin{pmatrix}2&1\\1&3\end{pmatrix}$, $B=\begin{pmatrix}1&0\\0&0\end{pmatrix}$, $C=\begin{pmatrix}0&0\\2&1\end{pmatrix}$.
\planteamiento calcular cada producto entrada a entrada.
\pasos
\begin{enumerate}[label=\arabic*.]
  \item $AB=\begin{pmatrix}2&0\\1&0\end{pmatrix}$ y $AC=\begin{pmatrix}2\cdot 0+1\cdot 2&2\cdot 0+1\cdot 1\\1\cdot 0+3\cdot 2&1\cdot 0+3\cdot 1\end{pmatrix}=\begin{pmatrix}2&1\\6&3\end{pmatrix}$.
  \item $AB+AC=\begin{pmatrix}4&1\\7&3\end{pmatrix}$.
  \item $B+C=\begin{pmatrix}1&0\\2&1\end{pmatrix}$ y $A(B+C)=\begin{pmatrix}2+2&0+1\\1+6&0+3\end{pmatrix}=\begin{pmatrix}4&1\\7&3\end{pmatrix}$.
  \item $BA=\begin{pmatrix}2&1\\0&0\end{pmatrix}$ y $CA=\begin{pmatrix}0&0\\2\cdot 2+1\cdot 1&2\cdot 1+1\cdot 3\end{pmatrix}=\begin{pmatrix}0&0\\5&5\end{pmatrix}$.
  \item $BA+CA=\begin{pmatrix}2&1\\5&5\end{pmatrix}$.
\end{enumerate}
\resultado \textbf{$AB+AC=A(B+C)=\begin{pmatrix}4&1\\7&3\end{pmatrix}$, mientras que $BA+CA=(B+C)A=\begin{pmatrix}2&1\\5&5\end{pmatrix}$ es distinta.} Cada distributiva vale por su lado, pero no se puede cambiar el lado.
\comprobacion $(B+C)A=\begin{pmatrix}1&0\\2&1\end{pmatrix}\begin{pmatrix}2&1\\1&3\end{pmatrix}=\begin{pmatrix}2&1\\4+1&2+3\end{pmatrix}=\begin{pmatrix}2&1\\5&5\end{pmatrix}$, que coincide con el paso~5.
\end{solucion}

\begin{solucion}{eje:4-coste}
\datos $A\in\R^{3\times 5}$, $B\in\R^{5\times 2}$, $C\in\R^{2\times 4}$; coste $mnp$ (\cref{prop:coste}).
\planteamiento sumar el coste de los dos productos de cada orden.
\pasos
\begin{enumerate}[label=\arabic*.]
  \item $(AB)C$: $AB$ cuesta $3\cdot 5\cdot 2=30$ y es $3\times 2$; $(AB)C$ cuesta $3\cdot 2\cdot 4=24$. Total: $54$.
  \item $A(BC)$: $BC$ cuesta $5\cdot 2\cdot 4=40$ y es $5\times 4$; $A(BC)$ cuesta $3\cdot 5\cdot 4=60$. Total: $100$.
  \item (b) $(u^{\top}v)w^{\top}$: el escalar $u^{\top}v$ cuesta $N$ y multiplicarlo por $w^{\top}$ cuesta $N$. Total: $2N$.
  \item $u^{\top}(vw^{\top})$: $vw^{\top}$ es $N\times N$ y cuesta $N^{2}$; $u^{\top}(vw^{\top})$ cuesta otras $N^{2}$. Total: $2N^{2}$.
  \item Con $N=4096$: $2N=8192$ y $2N^{2}=33\,554\,432$, un factor $N=4096$ mayor.
\end{enumerate}
\resultado \textbf{$(AB)C$ cuesta 54 multiplicaciones y $A(BC)$ cuesta 100. $(u^{\top}v)w^{\top}$ cuesta $2N$ y $u^{\top}(vw^{\top})$ cuesta $2N^{2}$.} El resultado es el mismo por la asociativa, pero el coste no.
\comprobacion ambos órdenes de (a) producen una matriz $3\times 4$; en (b), ambos producen el vector fila $(u^{\top}v)\,w^{\top}$ de $N$ componentes.
\end{solucion}

\begin{solucion}{eje:4-proyeccion}
\datos $A=\begin{pmatrix}1&0\\0&1\\0&0\end{pmatrix}$, $b=(4,-1,2)$, $a=(1,1,1)$.
\planteamiento calcular $P=AA^{\top}$ (filas de $A$ por filas de $A$) y aplicar $P^{2}=P$; en (c), $P_a=aa^{\top}/(a^{\top}a)$.
\pasos
\begin{enumerate}[label=\arabic*.]
  \item (a) $(AA^{\top})_{ij}$ es el producto escalar de las filas $i$ y $j$ de $A$, que son $(1,0)$, $(0,1)$ y $(0,0)$: $P=\diag(1,1,0)$.
  \item $Pb=(4,-1,0)$: se conservan las dos primeras coordenadas y se anula la tercera.
  \item (b) $P^{2}=\diag(1^{2},1^{2},0^{2})=P$.
  \item $(I-P)^{2}=I-2P+P^{2}=I-2P+P=I-P$, usando $P^{2}=P$.
  \item $I-P=\diag(0,0,1)$ y $(I-P)b=(0,0,2)$: proyecta sobre el eje $z$, perpendicular al plano generado por las columnas de $A$.
  \item (c) $a^{\top}a=3$ y $aa^{\top}$ es la matriz $3\times 3$ de unos; $P_a=\tfrac13 aa^{\top}$.
  \item $P_ab=\tfrac13(4-1+2)(1,1,1)=\tfrac53(1,1,1)\approx(1{,}667;\,1{,}667;\,1{,}667)$.
\end{enumerate}
\resultado \textbf{$P=\diag(1,1,0)$ y $Pb=(4,-1,0)$; $P^{2}=P$ y $(I-P)^{2}=I-P$; $I-P$ proyecta sobre el eje $z$; $P_ab=\tfrac53(1,1,1)$, como en el \cref{ej:proyeccion}.}
\comprobacion $Pb+(I-P)b=(4,-1,0)+(0,0,2)=b$ y $(4,-1,0)\cdot(0,0,2)=0$: $b$ se descompone en dos piezas ortogonales.
\end{solucion}

\begin{solucion}{eje:4-ortogonal}
\datos $R_\theta=\begin{pmatrix}\cos\theta&-\sin\theta\\ \sin\theta&\cos\theta\end{pmatrix}$; fórmulas del coseno y del seno de la suma.
\planteamiento calcular los productos entrada a entrada; para (c), usar $(AB)^{\top}=B^{\top}A^{\top}$.
\pasos
\begin{enumerate}[label=\arabic*.]
  \item (a) Las columnas de $R_\theta$ son $c_1=(\cos\theta,\sin\theta)$ y $c_2=(-\sin\theta,\cos\theta)$. $(R_\theta^{\top}R_\theta)_{ij}=c_i\cdot c_j$: $c_1\cdot c_1=c_2\cdot c_2=\cos^{2}\theta+\sin^{2}\theta=1$ y $c_1\cdot c_2=-\cos\theta\sin\theta+\sin\theta\cos\theta=0$. Luego $R_\theta^{\top}R_\theta=I_2$.
  \item (b) La primera columna de $R_\theta R_\varphi$ es
  \[
  R_\theta\begin{pmatrix}\cos\varphi\\ \sin\varphi\end{pmatrix}=\begin{pmatrix}\cos\theta\cos\varphi-\sin\theta\sin\varphi\\ \sin\theta\cos\varphi+\cos\theta\sin\varphi\end{pmatrix}=\begin{pmatrix}\cos(\theta+\varphi)\\ \sin(\theta+\varphi)\end{pmatrix}.
  \]
  \item La segunda columna es
  \[
  R_\theta\begin{pmatrix}-\sin\varphi\\ \cos\varphi\end{pmatrix}=\begin{pmatrix}-\cos\theta\sin\varphi-\sin\theta\cos\varphi\\ -\sin\theta\sin\varphi+\cos\theta\cos\varphi\end{pmatrix}=\begin{pmatrix}-\sin(\theta+\varphi)\\ \cos(\theta+\varphi)\end{pmatrix}.
  \]
  Por tanto, $R_\theta R_\varphi=R_{\theta+\varphi}$.
  \item Como $\theta+\varphi=\varphi+\theta$, $R_\theta R_\varphi=R_{\theta+\varphi}=R_\varphi R_\theta$.
  \item (c) Si $Q_1^{\top}Q_1=Q_2^{\top}Q_2=I$: $(Q_1Q_2)^{\top}(Q_1Q_2)=Q_2^{\top}Q_1^{\top}Q_1Q_2=Q_2^{\top}Q_2=I$.
  \item (d) Para $D=\diag(\lambda_1,\lambda_2)$ y $R=R_{90^\circ}$: $RD=\begin{pmatrix}0&-\lambda_2\\ \lambda_1&0\end{pmatrix}$ y $DR=\begin{pmatrix}0&-\lambda_1\\ \lambda_2&0\end{pmatrix}$, que coinciden si y solo si $\lambda_1=\lambda_2$.
\end{enumerate}
\resultado \textbf{Las rotaciones son ortogonales, se componen sumando ángulos y conmutan entre sí; el producto de ortogonales es ortogonal.} $S$ no conmuta con $R$ porque escala los dos ejes por factores distintos ($2\neq 0{,}5$) y no es ortogonal.
\comprobacion con $\theta=\varphi=90^\circ$: $R_{90^\circ}^{2}=\begin{pmatrix}0&-1\\1&0\end{pmatrix}^{2}=\begin{pmatrix}-1&0\\0&-1\end{pmatrix}=R_{180^\circ}$; con los datos del paso~6, $\lambda_1=2$ y $\lambda_2=0{,}5$ reproducen $RS$ y $SR$ del ejercicio~\ref{eje:4-rs}.
\end{solucion}

% ---------------------------------------------------------------------
\section{Tema 5}

\begin{solucion}{eje:5-tamano}
\datos $K=50\,257$; $d=768$ o $d=1600$; 4 bytes por número.
\planteamiento $E$ tiene una fila por token y una columna por dimensión: $Kd$ parámetros.
\pasos
\begin{enumerate}[label=\arabic*.]
  \item Tamaño: $E\in\R^{50\,257\times 768}$.
  \item Parámetros: $50\,257\cdot 768=50\,257\cdot 700+50\,257\cdot 68=35\,179\,900+3\,417\,476=38\,597\,376$.
  \item Memoria: $38\,597\,376\cdot 4=154\,389\,504$ bytes $\approx 154$~MB.
  \item Con $d=1600$: $50\,257\cdot 1600=80\,411\,200$ parámetros y $80\,411\,200\cdot 4=321\,644\,800$ bytes $\approx 322$~MB.
\end{enumerate}
\resultado \textbf{$E$ es $50\,257\times 768$, con $38\,597\,376$ parámetros ($\approx 154$~MB); con $d=1600$, $80\,411\,200$ parámetros ($\approx 322$~MB).}
\comprobacion orden de magnitud: $5\cdot 10^{4}\cdot 8\cdot 10^{2}=4\cdot 10^{7}\approx 3{,}9\cdot 10^{7}$.
\end{solucion}

\begin{solucion}{eje:5-onehot}
\datos $e_i\in\R^{K}$ one-hot; $E\in\R^{K\times d}$; en GPT-2, $K=50\,257$ y $d=768$.
\planteamiento escribir la componente $j$ del producto con la \cref{def:producto}.
\pasos
\begin{enumerate}[label=\arabic*.]
  \item $(e_i^{\top}E)_j=\sum_{k}(e_i)_kE_{kj}$.
  \item $(e_i)_k=0$ si $k\neq i$ y $(e_i)_i=1$, luego solo sobrevive el término $k=i$: $(e_i^{\top}E)_j=E_{ij}$.
  \item Esto vale para $j=1,\dots,d$: $e_i^{\top}E=(E_{i1},\dots,E_{id})$, la fila $i$.
  \item El producto requiere $Kd$ multiplicaciones; la lectura copia $d$ números. El cociente es $Kd/d=K=50\,257$.
\end{enumerate}
\resultado \textbf{$e_i^{\top}E$ es la fila $i$ de $E$; el producto requiere unas $50\,000$ veces más operaciones que la lectura directa.}
\comprobacion en el \cref{ej:embedding}, $(0,0,1,0)\,E=(0{,}5;\,0{,}5;\,0{,}5)$, que es la fila~2.
\end{solucion}

\begin{solucion}{eje:5-busqueda}
\datos $\hat E$ con filas unitarias; $K=50\,257$, $d=768$; lote de $B$ consultas.
\planteamiento $\cos(v_w,q)=\hat v_w\cdot\hat q$ y la lectura por filas del producto matriz-vector (\cref{prop:filas-columnas}).
\pasos
\begin{enumerate}[label=\arabic*.]
  \item (a) La componente $w$ de $\hat E\hat q$ es $\hat E_{w,:}\cdot\hat q=\cos(v_w,q)$. Un único producto matriz-vector da los $K$ cosenos. Como dividir por $\norm{q}>0$ no altera el orden (\cref{prop:coseno}(i)), basta calcular $\hat Eq$.
  \item (b) $\hat Eq$ requiere $Kd=50\,257\cdot 768=38\,597\,376$ multiplicaciones. La normalización de $E$ se hace una sola vez y no se repite en cada consulta.
  \item (c) Se apilan las consultas como columnas de $Q\in\R^{d\times B}$ y se calcula $\hat EQ\in\R^{K\times B}$, con $KdB$ multiplicaciones en un solo producto de matrices, que la GPU ejecuta eficientemente.
  \item Para vocabularios o colecciones mucho mayores se usan índices de búsqueda aproximada de vecinos, que evitan comparar con todos los candidatos (módulo~22).
\end{enumerate}
\resultado \textbf{Con filas normalizadas de antemano, los $K$ cosenos son $\hat Eq$ salvo un factor positivo común; cuesta unos $3{,}86\cdot 10^{7}$ multiplicaciones por consulta, y un lote se procesa con el producto $\hat EQ$.}
\comprobacion el coste coincide con el del producto $e_i^{\top}E$ del ejercicio~\ref{eje:5-onehot}, porque ambos recorren una vez todas las entradas de una matriz $K\times d$.
\end{solucion}

\begin{solucion}{eje:5-analogias}
\datos vectores del \cref{ej:juguete} y los cinco vectores geográficos del enunciado.
\planteamiento aritmética componente a componente y similitud coseno.
\pasos
\begin{enumerate}[label=\arabic*.]
  \item (a) $v_{\text{hombre}}\cdot v_{\text{mujer}}=0+(-1)+4+0=3$ y ambas normas valen $\sqrt5$: $\cos=3/5=0{,}6$.
  \item (b) $v_{\text{rey}}-v_{\text{hombre}}=(2,0,0,0)$; sumando $v_{\text{mujer}}$: $(2,-1,2,0)=v_{\text{reina}}$.
  \item (c) $v_{\text{Madrid}}-v_{\text{España}}=(-1,1,0,0)$; sumando $v_{\text{Francia}}$: $q=(0,1,0,1)=v_{\text{París}}$, luego $\cos(q,v_{\text{París}})=1$.
  \item $q\cdot v_{\text{Lyon}}=0{,}5+1=1{,}5$; $\norm{q}=\sqrt2\approx 1{,}414$; $\norm{v_{\text{Lyon}}}=\sqrt{1{,}25}\approx 1{,}118$.
  \item $\cos(q,v_{\text{Lyon}})=1{,}5/\sqrt{2{,}5}=1{,}5/1{,}581\approx 0{,}9487$.
\end{enumerate}
\resultado \textbf{(a) $0{,}6$; (b) se obtiene exactamente «reina»; (c) $q=(0,1,0,1)$, con coseno~1 con París y $\approx 0{,}9487$ con Lyon.} París gana, pero Lyon queda cerca.
\comprobacion $0{,}9487<1$, como exige el \cref{teo:cs}; y (b) es el caso exacto de la \cref{prop:paralelogramo}.
\end{solucion}

\begin{solucion}{eje:5-secuencia}
\datos $E\in\R^{4\times 3}$ del \cref{ej:embedding}; secuencia $(2,0,2)$.
\planteamiento \cref{prop:onehot}: cada fila de $O$ es el one-hot de un token y $X=OE$.
\pasos
\begin{enumerate}[label=\arabic*.]
  \item (a) Columna 1: $0\cdot 0{,}1+0\cdot 1+1\cdot 0{,}5+0\cdot(-0{,}2)=0{,}5$. Las columnas 2 y 3 dan también $0{,}5$.
  \item (b) $O=\begin{pmatrix}0&0&1&0\\1&0&0&0\\0&0&1&0\end{pmatrix}\in\R^{3\times 4}$.
  \item $X=OE$: sus filas son las filas 2, 0 y 2 de $E$:
  \[
  X=\begin{pmatrix}0{,}5&0{,}5&0{,}5\\0{,}1&0{,}2&0\\0{,}5&0{,}5&0{,}5\end{pmatrix}\in\R^{3\times 3}.
  \]
\end{enumerate}
\resultado \textbf{(a) $e_2^{\top}E=(0{,}5;\,0{,}5;\,0{,}5)$. (b) $X$ es $3\times 3$ ($n\times d$), con las filas primera y tercera iguales.}
\comprobacion tamaños: $(3\times 4)(4\times 3)=3\times 3$.
\end{solucion}

\begin{solucion}{eje:5-logits}
\datos $E$ del \cref{ej:embedding}; $h=(1,0,-1)$.
\planteamiento $z_i=E_{i,:}\cdot h$ y $\softmax(z)_i=e^{z_i}/\sum_je^{z_j}$.
\pasos
\begin{enumerate}[label=\arabic*.]
  \item $z_0=0{,}1-0=0{,}1$; $z_1=1+1=2$; $z_2=0{,}5-0{,}5=0$; $z_3=-0{,}2-0{,}1=-0{,}3$.
  \item Exponenciales: $e^{0{,}1}=1{,}105$; $e^{2}=7{,}389$; $e^{0}=1$; $e^{-0{,}3}=0{,}7408$. Suma: $10{,}24$.
  \item Softmax: $1{,}105/10{,}24=0{,}1080$; $7{,}389/10{,}24=0{,}7219$; $1/10{,}24=0{,}0977$; $0{,}7408/10{,}24=0{,}0724$.
  \item La fila~1, $(1,0,-1)$, es igual a $h$: $\cos(E_{1,:},h)=1$ y $z_1=\norm{E_{1,:}}\norm{h}\cos=\sqrt2\cdot\sqrt2\cdot 1=2$.
\end{enumerate}
\resultado \textbf{$z=(0{,}1;\,2;\,0;\,-0{,}3)$ y $\softmax(z)\approx(0{,}1080;\,0{,}7219;\,0{,}0977;\,0{,}0724)$. El token más probable es el~1, cuyo embedding apunta exactamente en la dirección de $h$.}
\comprobacion $0{,}1080+0{,}7219+0{,}0977+0{,}0724=1{,}0000$.
\end{solucion}

\begin{solucion}{eje:5-atencion}
\datos (a) $q=(1,0)$, $k_1=(1,1)$, $k_2=(0,2)$, $k_3=(-1,0)$, $d_k=2$. (b) Componentes independientes de media~0 y varianza~1. (c) $d=64$.
\planteamiento (a) puntuaciones y softmax; (b) linealidad de la esperanza y aditividad de la varianza para sumandos independientes.
\pasos
\begin{enumerate}[label=\arabic*.]
  \item (a) $q\cdot k_1=1$, $q\cdot k_2=0$, $q\cdot k_3=-1$; divididos por $\sqrt2\approx 1{,}414$: $(0{,}7071;\,0;\,-0{,}7071)$.
  \item $e^{0{,}7071}=2{,}028$, $e^{0}=1$, $e^{-0{,}7071}=0{,}4931$; suma $3{,}521$; softmax $(0{,}5760;\,0{,}2840;\,0{,}1400)$.
  \item (b) $q\cdot k=\sum_{i=1}^{d}q_ik_i$. Por independencia, $\E[q_ik_i]=\E[q_i]\E[k_i]=0$, luego $\E[q\cdot k]=0$.
  \item $\Var(q_ik_i)=\E[q_i^{2}k_i^{2}]-0^{2}=\E[q_i^{2}]\,\E[k_i^{2}]=1\cdot 1=1$, porque $\E[X^{2}]=\Var(X)+\E[X]^{2}=1$.
  \item Los $d$ sumandos $q_ik_i$ dependen de pares disjuntos de variables independientes, luego son independientes y sus varianzas se suman: $\Var(q\cdot k)=d$.
  \item (c) Con $d=64$, la desviación típica es $\sqrt{64}=8$. La softmax de $(8,0,-8)$ es $(2981;\,1;\,0{,}000335)/2982\approx(0{,}99966;\,0{,}00034;\,1{,}1\cdot 10^{-7})$: casi un vector one-hot.
  \item La softmax de $(1,0,-1)$ es $(2{,}718;\,1;\,0{,}3679)/4{,}086\approx(0{,}6652;\,0{,}2447;\,0{,}0900)$: una distribución repartida.
\end{enumerate}
\resultado \textbf{(a) puntuaciones $(0{,}7071;\,0;\,-0{,}7071)$ y pesos $(0{,}576;\,0{,}284;\,0{,}140)$. (b) $\E[q\cdot k]=0$ y $\Var(q\cdot k)=d$. (c) Sin dividir por $\sqrt d$ la softmax se satura; dividiendo, la distribución queda repartida y los gradientes no se anulan.}
\comprobacion (a) $0{,}576+0{,}284+0{,}140=1{,}000$. (b) Una simulación con $10^{5}$ pares de vectores gaussianos de dimensión 64 da una varianza muestral de $63{,}8$ antes de dividir por~8 y de $1{,}00$ después.
\end{solucion}


\backmatter
% =====================================================================
\begin{thebibliography}{99}
\addcontentsline{toc}{chapter}{Bibliografía}
\markboth{Bibliografía}{}
\small

\bibitem{bhk}
A.~Blum, J.~Hopcroft y R.~Kannan.
\newblock \emph{Foundations of Data Science}.
\newblock Cambridge University Press, 2020.
\newblock Cap.~2, «High-Dimensional Space». Texto del curso CS~4850 de Cornell.
\newblock \url{https://www.cs.cornell.edu/jeh/book.pdf}

\bibitem{bolukbasi}
T.~Bolukbasi, K.-W. Chang, J.~Zou, V.~Saligrama y A.~Kalai.
\newblock Man is to computer programmer as woman is to homemaker? Debiasing word embeddings.
\newblock En \emph{Advances in Neural Information Processing Systems 29}, 2016.
\newblock \url{https://arxiv.org/abs/1607.06520}

\bibitem{vmls}
S.~Boyd y L.~Vandenberghe.
\newblock \emph{Introduction to Applied Linear Algebra: Vectors, Matrices, and Least Squares}.
\newblock Cambridge University Press, 2018.
\newblock Texto del curso EE103 de Stanford.
\newblock \url{https://web.stanford.edu/~boyd/vmls/}

\bibitem{brown}
T.~B. Brown \emph{et al.}
\newblock Language models are few-shot learners.
\newblock En \emph{Advances in Neural Information Processing Systems 33}, 2020.
\newblock \url{https://arxiv.org/abs/2005.14165}

\bibitem{mml}
M.~P. Deisenroth, A.~A. Faisal y C.~S. Ong.
\newblock \emph{Mathematics for Machine Learning}.
\newblock Cambridge University Press, 2020.
\newblock Cap.~2, «Linear Algebra», y cap.~3, «Analytic Geometry».
\newblock \url{https://mml-book.github.io/}

\bibitem{firth}
J.~R. Firth.
\newblock A synopsis of linguistic theory, 1930--1955.
\newblock En \emph{Studies in Linguistic Analysis}, págs.~1--32. Blackwell, Oxford, 1957.

\bibitem{harris}
Z.~S. Harris.
\newblock Distributional structure.
\newblock \emph{Word}, 10(2--3):146--162, 1954.

\bibitem{jm3}
D.~Jurafsky y J.~H. Martin.
\newblock \emph{Speech and Language Processing}, 3.ª ed., borrador del 19 de agosto de 2026.
\newblock Cap.~5, «Embeddings».
\newblock \url{https://web.stanford.edu/~jurafsky/slp3/}

\bibitem{mikolov-a}
T.~Mikolov, K.~Chen, G.~Corrado y J.~Dean.
\newblock Efficient estimation of word representations in vector space.
\newblock arXiv:1301.3781, 2013.
\newblock \url{https://arxiv.org/abs/1301.3781}

\bibitem{mikolov-b}
T.~Mikolov, W.-t. Yih y G.~Zweig.
\newblock Linguistic regularities in continuous space word representations.
\newblock En \emph{Proceedings of NAACL-HLT}, págs.~746--751, 2013.
\newblock \url{https://aclanthology.org/N13-1090/}

\bibitem{nissim}
M.~Nissim, R.~van Noord y R.~van der Goot.
\newblock Fair is better than sensational: man is to doctor as woman is to doctor.
\newblock \emph{Computational Linguistics}, 46(2):487--497, 2020.

\bibitem{press}
O.~Press y L.~Wolf.
\newblock Using the output embedding to improve language models.
\newblock En \emph{Proceedings of EACL}, vol.~2, págs.~157--163, 2017.
\newblock \url{https://aclanthology.org/E17-2025/}

\bibitem{radford}
A.~Radford, J.~Wu, R.~Child, D.~Luan, D.~Amodei e I.~Sutskever.
\newblock Language models are unsupervised multitask learners.
\newblock Informe técnico, OpenAI, 2019.

\bibitem{strang}
G.~Strang.
\newblock \emph{Introduction to Linear Algebra}, 4.ª ed.
\newblock Wellesley-Cambridge Press, 2009.
\newblock Texto de los cursos 18.06 y 18.06SC del MIT. Hay 5.ª ed. (2016) y 6.ª ed. (2023); la numeración de secciones y problemas citada es la de la 4.ª ed.

\bibitem{mit1806}
G.~Strang.
\newblock 18.06: Linear Algebra.
\newblock Massachusetts Institute of Technology, primavera de 2010. Hoja de problemas~1.
\newblock \url{https://web.mit.edu/18.06/www/Spring10/}

\bibitem{mit1806sc}
G.~Strang.
\newblock 18.06SC: Linear Algebra.
\newblock Massachusetts Institute of Technology, MIT OpenCourseWare, otoño de 2011.
\newblock \url{https://ocw.mit.edu/courses/18-06sc-linear-algebra-fall-2011/}

\bibitem{vaswani}
A.~Vaswani, N.~Shazeer, N.~Parmar, J.~Uszkoreit, L.~Jones, A.~N. Gomez, Ł.~Kaiser e I.~Polosukhin.
\newblock Attention is all you need.
\newblock En \emph{Advances in Neural Information Processing Systems 30}, 2017.
\newblock \url{https://arxiv.org/abs/1706.03762}

\bibitem{cs224n}
D.~Yang y Y.~Choi.
\newblock CS224N: Natural Language Processing with Deep Learning.
\newblock Stanford University, invierno de 2026. Clase~2, «Word Vectors».
\newblock \url{https://web.stanford.edu/class/cs224n/}

\bibitem{d2l}
A.~Zhang, Z.~C. Lipton, M.~Li y A.~J. Smola.
\newblock \emph{Dive into Deep Learning}, versión 1.0.3.
\newblock Cambridge University Press, 2023.
\newblock \url{https://d2l.ai/}

\end{thebibliography}


\end{document}
